\chapter{Groups of diffeomorphismes}
\label{grp.diff}
\section{Introduction}

Shapes, alone or within images, are typically subject to rigid or
non-rigid deformations. In order to compare two of them, or to extract common
features, one typically needs to develop a matching procedure, which
consists in finding a generally one-to-one correspondence between the shapes.

From a formal point of view, we will consider ``objects'' as functions
defined on some set $\Om$ with values in some set $F$.  For example,
parameterized curves will be defined on an interval, with values in
$\mR^2$. An image is defined on a subset of $mR^2$ (or $\mR^3$), with
values in $\mR$. Similar objects will have similar features (corners,
patterns), at similar, but not identical, locations. The matching
problem consists in aligning the objects, which is a one-to-one map
from $\Om$ onto $\Om$ that register the features between the two objects.

To simplify, we will assume that $\Om$ is a bounded subset of $\mR^n$
(although the theory can be extended to unbounded sets). Its boundary
will be denoted  $\prt\Om$, and its closure $\bOm$. Deformable objects
are functions $m: \bOm \mapsto F$, with $F\sub \mR^d$. Matching two
objects $m$ and $\tm$ consists in defining a bijection $g: \bOm\mapsto
\bOm$ wihich is not too far from the identity, and such that
$m\circ g$ is close to $\tm$.

The issue, of course, is to rigorously define these measures of
proximity, between bijections and objects. In this part, are goals are
as follows.

\begin{itemize}
\item Build a group $G$ that contains diffeomorphisms of $\Om$
\item Equip  $G$ with a differential structure that allows for the
composition of infinitesimal elements, like for Lie groups.
\item Equip $G$ with a metric associated to a Riemannian structure
(\ref{geo.diff}) to evaluate the size of deformations.
\end{itemize}
\index{diffeomorphism}

We will not build here a Lie group of diffeomorphism (Lie groups in
infinite dimensions generally are intricate objects; for
diffeomorphisms, they generally require using $C^\infty$
functions). One of the simple way to define a group of diffeomorphisms
is by
$$
G_k = \defset{g: C^k \text{-diffeomorphisme sur } \overline{\Om}, g^{-1} \text{
    de classe } C^k}.
$$
Although this group can be equipped with a structure of complete
metric space
\cite{mic72, df01}, it is not suitable for differential calculs. One
practical way do induce infinitesimal perturbations in $G_k$ is via
the action of flows of differentla equations  (\cite{df01}). But these
flows can be used directly to build new groups of diffeomorphisms,
typically different from  $G_k$ \cite{tro98}. We will follow this
construction here.


As a motivation for the inroduction of flows, let us study the idea of
building diffeomorphisms by composing small deformations. Let a
diffeomorphism near the identity be written as $h(x) = x + \de v(x)$,
$\de$ being an arbitrarily small number. It is not hard to show that
this is a diffeomorphism for $\de \leq \de_0$, $\de_0$ depending on
$\norm{v}_\infty$ and $\norm{Dv}_\infty$. Consider now $n$ such
deformations,  $h_i = \id +\de v_i$,
$i=1, \ldots, n$, and denote $g_n = h_n\circ\cdots\circ h_1$. We have
the formula
$$
g_{n+1} = (\id + \de v_n) \circ g_n = g_n + \de v_n\circ g_n.
$$
Write this under the form
$$
\frac{1}{\de} (g_{n+1} - g_n) = g_n\circ v_n;
$$
for small $\de$, this equation can be interpreted as a time
discretization of an ordinary differential equation of the form
\begin{equation}
\label{flot.1}
\der{\bfg}{t} = \bfv(t, \bfg).
\end{equation}
Here, $\bfv$ is a time dependent vector field on $\Om$. It is the main
ingredient in the construction of groups of diffeomorphism proposed in
 \cite{tro98} (see also
\cite{dgm98}).

\section{Flows of diffeomorphismes}
\index{flows of diffeomorphisms}
\subsection{Construction}

The general principle is to fix a suitable functional space to which
the time-dependent vector field $\bfv$ belongs, and to define the
elements of $H$ to be solutions of  (\ref{flot.1}). It has to be
ensured that, for the considered  $\bfv$, (\ref{flot.1}) has solutions
over all times, and that they form a group. This is the goal of the
following definitions.

Let $L^\infty(\bOm, \mR^n)$ be the set of bounded measurable functions
from $\bOm\sub \mR^n$ to $\mR^n$. Let $C_0^1(\Om,
\mR^n)$ be the set of all $C^1$ vector fields $v$ on  $\bOm$ such that
$v$ and $Dv$ vanish on  $\prt\Om$. We will use the norm, for which
$C_0^1(\Om, \mR^n)$ is a Banach space:
$$\|v\|_{1, \infty} = \|v\|_\infty + \|dv\|_\infty\,.$$

We define the set $\CW$ of trajectories $\bfv: [0, 1] \ria
C_0^{1}(\Om,\mR^n)$ for which
$$
\|\bfv\| := \int_0^1 \|\bfv(t)\|^2_{1, \infty} dt
$$
is finite.
For all $t$, $\bfv(t)$ is a vector field on $\Om$, that we shall also
denote $\bfv(t, .)$. We also let $\CW^\infty$ be the subset of  $\CW$
that contains time-dependent vector fields  $\bfv$ such that $\bfv(t, .)\in
C^\infty(\Om, \mR^n)$ at all times.


For $\bfv\in\CW$, and for $x\in \bOm$, denote $\bfg_{\bfv}(.,
x)$ the solution (if it exists) of the ODE (\ref{flot.1}) with initial
condition $\bfv(0, x) = x$.

\encadre{
\begin{theorem}
\label{flot.exist}
If $\bfv\in\CW$, then, for all $x\in \bOm$, $\bfg_{\bfv}(t, x)$ exists
for all 
$t\in[0,1]$ and is uniquely defined. At all time $t$, the function $x\mapsto
\bfg_\bfv(t, x)$ is a diffeomorphism of $\Om$. Its differential is the
solution of the linear differential equation
\begin{equation}
\label{flot.diff}
\pdr{}{t}d_x\bfg_\bfv(t, .) = d_{\bfg_\bfv(t, x)} \bfv(t,
.).d_x\bfg_\bfv(t, .)
\end{equation}
with initial condition $d_x\bfg_\bfv(0, .) = \Id$. 
\end{theorem}}
\begin{proof}
We first show existence and uniqueness. Fix $t\in [0, 1]$ and $\de
>0$. Let $I = I(t, \de)$ denote the interval
$[0,1]\,\cap\, ]t-\de, t+\de[$. Fix also $x\in\bOm$.

If $\phi$ is a continuous function from $I$ to $\mR^3$ such that 
$\phi(t)=x$, define, for $s\in I$
$$
\Ga(\phi)(s) = x + \int_t^s \bfv(u, \phi(u)) du
$$
The function $\Ga(\phi)$ is also continuous on $I$ and is such that 
$\Ga(\phi)(t) = x$. The set of all such function being a Banach space
for the supremum norm, it suffices to show that $\Ga$ is contractant
for small enough $\de$ to prove the existence of a unique fixed point.

From Schwartz inequality, we have
%\begin{multline}
\begin{equation}
\label{eq:***} \|\Ga(\phi) - \Ga(\psi)\|_\infty \leq
\|\phi-\psi\|_\infty \int_{s\wedge t}^{s\vee t} \|\bfv(u, .)\|_{1,
\infty} du
%\\ 
\leq
\sqrt{|t-s|}\,\|\bfv\|\,\|\phi-\psi\|_\infty.
\end{equation}
%\end{multline}
So its suffices to take  $\de < \|\bfv\|^{-2}$ to ensure that $\Ga$ is
contractant. This shows that a unique solution (given its value at
$t$) of  (\ref{flot.1})
exists on intervals of radius $\de$ around $t$. But we have found that
$\de$ can be chosen uniformly in $t$, which implies that the solution
exists and is unique over all $[0,1]$ (and in fact any time interval),
since one can extend any solution by intervals of radii smaller than
$\de$.

The continuity of  $\bfg_\bfv$ in its initial conditions is a standard
consequence of Gronwall's lemma that we give without proof:
\begin{theorem}[Gronwall's lemma]
\label{th:gronwall}
Assume that $\phi, \psi$ and $w$ and bounded, positive functions on an
interval $[0, c]$, and that
$$
w(t) \leq \phi(t) + \int_0^t \psi(s) w(s) ds.
$$
Then
$$
w(t) \leq \phi(t) + \int_0^t \phi(s)\psi(s) e^{\int_s^t \psi(u)du}ds.
$$
\end{theorem}

 

If $x,
y\in\bOm$, we have
\begin{eqnarray*}
|\bfg_\bfv(t, x) - \bfg_\bfv(t, y)| &=& \left| x-y + \int_0^t
\left(\bfv(s, \bfg_\bfv(s, x)) - \bfv(s, \bfg_\bfv(s,
y))\right)ds\right|\\
&\leq & |x-y| + \int_0^t \|\bfv(s, .)\|_{1, \infty} |\bfg_\bfv(s, x) -
\bfg_\bfv(s, y)| ds\\
&\leq& |x-y| + \sqrt{\int_0^t \|\bfv(s, .)\|^2_{1, \infty}ds \int_0^t |\bfg_\bfv(s, x) -
\bfg_\bfv(s, y)|^2 ds}
\end{eqnarray*}
Letting $M_t = |\bfg_\bfv(s, x) - \bfg_\bfv(s, y)|^2$,
we get
$$
M_t \leq 2|x-y|^2 + 2\|\bfv\|^2 \int_0^t M_s ds
$$
which implies
\begin{equation}
\label{asc.1}
M_t \leq 2|x-y| \exp(2\|\bfv\|^2),
\end{equation}
which proves the continuity of  $\bfg_\bfv$ with respect to $x$.

To show that $\bfg_\bfv(t, .)$ is a homeomorphism (i.e., invertible
with a continuous inverse), it suffices to notice that, for all $t>0$,
the function $\bfw(s, .) = -\bfv(t-s,
.)$ for  $s\in [0, t]$ and $\bfw(s, .) = 0$ pour $s > t$ belongs to
$\CW$, and is such that $\bfg_\bfw(t,.) =
[\bfg_\bfv(t, .) ]^{-1}$. Indeed, define $\bfh : (s, x)\mapsto
\bfg_\bfv(t-s, x)$. We have $bfh(s,x) = bfg_\bfw(s, \bfg_\bfv(t, x))$,
because these two functions are solutions of 
$$
\pdr{\psi}{s}(s,x) = - \bfv(t-s, \psi(s, x))
$$
with initial conditions $\psi(0, x) = \bfg(t, x)$. Therefore, 
$$
\bfg_\bfw(t, \bfg_\bfv(t, x)) = \bfh(1, x) = \bfg_\bfv(0, x) = x$$
which shows that $\bfg_\bfv(t, .)$ is invertible, with a continuous inverse
(since it also comes from an element of $\CW$), hence is a
homeomorphism of $\bOm$.

Let's show now that  $\bfg_\bfv(t, .)$ is a diffeomorphism. Assuming
that this is the case, we can formally compute the differential of the
equation
$$
\pdr{\bfg_\bfv}{t}(t, x+\ep\de) = \bfv(t, \bfg_\bfv(t, x+\ep\de))
$$ at $\ep=0$ and get
$$
\pdr{}{t}d_x\bfg_\bfv(t,.).\de = d_{\bfg_\bfv(t,x)}\bfv(t, .)d_x\bfg_\bfv(t,.).\de.
$$
In view of this,, let's introduce the linear differential equation
$$
\pdr{W}{t} = d_{\bfg_\bfv(t,x)}\bfv(t, .) W
$$
with initial condition $W(0) = \de$. This equation has a unique
solution on any interval (by the same kind of argument used to prove
that \eqref{flot.1} has a unique solution). The previous formal
computation shows that this solution should coincide with
$d_x\bfg_\bfv(t, .).\de$, but this remains to prove.

For this, introduce
$$
a_\ep(t) = \left(\bfg_\bfv(t, x+\ep\de) - \bfg_\bfv(t, x)\right)/\ep -
W_t.
$$
We need to prove that $a_\ep(t)$ tends to 0 when $\ep$ tends to
0. Introduce a modulus of continuity for $d_x\bfv(t, .)$ on
$\Om$, letting
$$
\mu_t(\al) = \max\defset{\abs{d_x\bfv(t, .) - d_y\bfv(t, .)}: x,
  y\in\Om, \abs{x-y} \leq \al}.
$$
Since $d_x\bfv(t, .)$ is uniformly continuous on the compact set 
$\overline{\Om}$, we know that $\mu_t(\al)\to 0$ if $\al\to 0$. We can
have, using the intehral version of the differential equations that
define  $\bfg_\bfv$
and $W$: 
\begin{align*}
a_\ep(t) &= \frac{1}{\ep}\int_0^t \left(\bfv(s, \bfg_\bfv(s, x+\ep\de)) - \bfv(s,
\bfg_\bfv(s, x))\right) ds - \int_0^t   d_{\bfg_\bfv(s,x)}\bfv(s,.) W(s) ds \\
&= \int_0^t d_{\bfg_\bfv(s,x)}\bfv(s,.) a_\ep(s) ds \\
&+
\frac{1}{\ep}\int_0^t \left(\bfv(s, \bfg_\bfv(s, x+\ep\de)) - \bfv(s,
\bfg_\bfv(s, x))\right. \\
&\ \ \ \ \ \ -\left. \ep d_{\bfg_\bfv(s,x)}\bfv(s,.) \left(\bfg_\bfv(s,
x+\ep\de) - \bfg_\bfv(s, x)\right) \right) ds
\end{align*}
Using the fact that, for all $x, y\in\Om$:
$$
\abs{\bfv(t, y) - \bfv(t, x) - d_x\bfv(t, .)(y-x)} \leq \mu_s(\abs{x-y})
\abs{x-y}
$$
and combining this to \eqref{asc.1}, we can write 
$$
\abs{a_\ep(t)} \leq  \int_0^t \norm{\bfv(s, .)}_{1, \infty}
\abs{a_\ep(s)} ds + C(\bfv) \abs{\de} \int_0^1 \mu_s(\ep
C(\bfv)\abs{\de}) ds
$$
for some constant $C(\bfv)$ which only depends on $\bfv$. The
convergence of $a_\ep$ to 0 will be obtained as a consequence of
Gromwall's lemma as soon as we will have proved that 
$$
\lim_{\al \to 0}  \int_0^1 \mu_s(\al)ds = 0.
$$
But this fact is a consequence of the pointwise convergence of $\mu_s$
to 0, of the upper bound $\mu_s(\al) \leq 2 \norm{\bfv}_{1, \infty}$
and of the dominated convergence theorem.
\end{proof}

The inequality (\ref{asc.1}) shows that, if $\CV$ is a bounded subset
of  $\CW$, the family of functions $\bfg_\bfv(t, .)$ for $\bfv\in \CV$
is equicontinuous (i.e., has a common modulus of continuity). The same
tachnique can be used to prove the analogous result for  $\bfg_\bfv(.,
x)$, and therefore for the family $\bfg_\bfv$ of time dependent vector
fileds. Ascoli's theorem therefore implies that the map $\bfv \mapsto
\bfg_\bfv$ is compact from 
$\CW$ to $C^0([0,1]\ti\bOm, \bOm)$ (i.e., carries bounded sets onto
compact sets).

The other fundamental result that can be deduced from the previous
proof is that the image of  $\CW$ by the map $\bfv
\mapsto \bfg_\bfv$ is a subgroup of the group of diffeomorphisms of
$\bOm$. Indeed, the stability by composition is obtained by
concatenating two flows, and the existence of the inverse has been
shown in the proof of theorem
\ref{flot.exist}. We can therefore provide the following definition.
\begin{definition}
\label{g.infty}
We let $G = G(\CW)$ be the group of diffeomorphisms of $\bOm$ defined by
$$
G = \{\bfg_\bfv(1, .), \bfv\in\CW\}
$$
\end{definition}
(There is no loss of generality in restriction to diffeomorphisms at
time 1since $\bfg_{t\bfv}(1, .) = \bfg_\bfv(t,, .)$.)

The group which is defined in this way is consistent with the
intuition of constructing diffeomorphisms by iterating infinitesimal
deformations. This approach also provides a distance on $G$, letting,
for  $g, g'\in G$
$$
d_\CW(g, g')^2 = \inf\defset{\int_0^1 \|\bfv(t, .)\|_{1, \infty}^2
dt: \bfg_\bfv(1, g(x)) \equiv g'(x)\|}
$$

The fact that $d$ is a distance will be proved later in a more general
case. Another important  property of the space  $\CW$ is the behavior
of the solutions of  (\ref{flot.1}) when $\bfv$ weakly converges in
the following sense.

\begin{definition}
\label{weak*}
A sequence $\bfv_n$ in $\CW$ weakly converges to $\bfv\in \CW$ if and
only ifm for all $C^\infty$ function, $\phi:[0,1]\ti\overline{\Om} \ria
\mR$, 
we have 
$$
\int_0^1\int_\Om \phi(t, x) \bfv_n(t, x)dt dx \ria \int_0^1\int_\Om
\phi(t, x) \bfv(t, x)dt dx
$$
\end{definition}

We then have the following theorem
\medskip
\encadre{
\begin{theorem}
\label{weak.closed}
If $\bfv_n$ is a bounded sequence in $\CW$ that weakly converges to
$\bfv^*\in \CW$, then $\bfg_{\bfv_n}$ converges uniformly on $[0,
1]\ti\bOm$ to $\bfg_{\bfv^*}$.
\end{theorem}}

\begin{proof}
The proof here is adapted from \cite{dgm98}. Since we know that $\bfv \mapsto
\bfg_\bfv$ is compact, it suffices to show that any subsequence of
$\bfg_{\bfv_n}$ that converges uniformly also converges to 
$\bfg_{\bfv^*}$. Fix such a subsequence, that will still be denoted
$\bfv_n$, and let  $\bfg$ be its
limit. To show that $\bfg = \bfg_{\bfv^*}$, it suffices to prove that,
for all $t$ and all $x$
$$
\bfg(t, x) = x + \int_0^t \bfv^*(s, \bfg(s, x)) ds.
$$
Since this identity is true for  $\bfg_{\bfv_n}$ (with $\bfv_n$ in
place of  $\bfv^*$), and since $\bfg_{\bfv_n}$ converges uniformly to
$\bfg$, all we have to show is that, for all $t$ and $x$,
$$
\int_0^t [\bfv^*(s, \bfg(s, x)) - \bfv_n(s, \bfg_{\bfv_n}(s, x))] ds
\ria 0.
$$
We have
\begin{multline*}
\left|\int_0^t [\bfv^*(s, \bfg(s, x)) - \bfv_n(s, \bfg_{\bfv_n}(s,
x))] ds\right|
\leq \\
\int_0^t
\left|\bfv_n(s, \bfg(s, x)) - \bfv_n(s, \bfg_{\bfv_n}(s,
x))\right| ds
+ \left|\int_0^t \bfv^*(s, \bfg(s, x)) - \bfv_n(s,
\bfg(s,x)) ds \right|
\end{multline*}
For the first integral, write
\begin{eqnarray*}
\int_0^t
\left|\bfv_n(s, \bfg(s, x)) - \bfv_n(s, \bfg_{\bfv_n}(s,
x))\right| ds &\leq& \int_0^1 \|\bfv_n(t, .)\|_{1, \infty} |\bfg(s, x)
- \bfg_{\bfv_n}(s, x)| ds\\
&\leq& \|\bfg-\bfg_{\bfv_n}\|_\infty \|\bfv_n\|.
\end{eqnarray*}
Since $\|\bfv_n\|$ is bounded and 
$\bfg_{\bfv_n}$ uniformly converges to $\bfg$, this integral converges
to 0. For the second integral, we can use the following lemma:
\begin{lemma}
\label{lemme.weak}
If $\bfv_n$ is bounded in $\CW$ and weakly converges to $\bfv^*$,
then, for all continuous function $\phi: [0, 1]\to \bOm$, we have, for
all
$t\in [0,1]$,
\begin{equation}
\label{eq.weak}
\lim_{n\ria\infty} \int_0^t (\bfv_n(s, \phi(s)) - \bfv^*(s, \phi(s)))
ds  = 0.
\end{equation}
\end{lemma}
Since any continuous function can be uniformly approached by a piecewise
constant function, it suffices to prove this for a piecewise constant
function (we skip the limiting argument here, which can be done like
we have estimated the first integral, using the fact that $\bfv_n$ is
bounded in $\CW$). Moreover,  as piecewise constant
functions, we obviously only need to consider  $\phi$ of the form $\phi(s) =
x.\bfone_A(s)$ where $A= [t_1, t_2]$ is a subinterval of  $[0,
1]$. Writing the associated integral as a difference between an
integral on $[0,
t_1]$ and an integral on $[0, t_2]$, we see that the lemma only needs
to be proved for all $t$ and constant $\phi$. So let's assume that
$\phi(s) = x$ and denote 
$$
Q_n(t, x) = \int_0^t (\bfv_n(s, x) - \bfv^*(s, x)) ds.
$$
Proceding like for equation (\ref{eq:***}), we show that 
$Q_n$ is an equicontinuous family of functions $[0,
1]\ti\bOm$; since $Q_n(0, x) = 0$ for all $x$, Azrela-Ascoli's theorem
implies that this sequence is relatively compact. It therefore
suffices to show that any converging subsequence of  $Q_n$ converges to 0.
Fix such a subsequence, still denoting it  $Q_n$ to save on notation,
and let $Q$ be its limit. Let $f(t, x)$ be a $C^\infty$ function on
$[0, 1]\ti \bOm$. We have 
$$
\int_0^1\int_\Om Q^n(t, x) f(t, x) dt dx = \int_0^1 \int_\Om
(\bfv_n(s, x) - \bfv^*(s, x)) \left(\int_s^1 f(t, x) dt\right)ds dx.
$$
From the weak convergence of  $\bfv_n$, the right-hand term tends to
0. Since  $Q_n$ converges uniformly to $Q$, we have
$$
\int_0^1\int_\Om Q(t, x) f(t, x) dt dx = 0
$$ 
for all $C^\infty$ $f$, which implies the desired result.
\end{proof}

%%%%%%%%%%%%%%%%%%%%%%%%%%%%%%%%%%

\subsection{Extension}
\label{ext.diff}

The group $G$ in the previous section has been built using vector
fields in the space $\CW = C^{1}_0(\Om, \mR^n)$. Clearly, all the properties
that have been shown above remain valid in this space is replaced by a
smaller Banach spaces (i.e., that can be continuously embedded in
it). Following  \cite{tro98}, we introduce the definition
\encadre{
\begin{definition}
\label{adm}
A Banach space $\CB\subset L^2(\Om, \mR^n)$, with norm $\|.\|_\CB$,
is admissible if and only if $C^\infty_0(\Om, \mR^n)$ is dense in
$\CB$ and there exists a constant $K > 0$ such that, for all
$v\in C^\infty_0(\Om, \mR^n)$, we have
$$
\|v\|_{1, \infty} \leq K\|v\|_\CB
$$
\end{definition}}

All the results that precede remain true if the norm $\|\,.\,\|_{1,
\infty}$ is replaced by $\|\,.\,\|_\CB$. One can therefore build a
group $G(\CB) \sub G(\CW)$ and a new distance  $d_\CB$ on $G(\CB)$
defined by
$$
d_\CB(g, g')^2 = \inf\defset{ \int_0^1 \|\bfv(t, .)\|_\CB^2 dt:
\bfv(1, g(x)) \equiv g'(x)\|}
$$

Using Sobolev inclusions \cite{ada75,bre83}, $\CB$ can be chosen to be
some $W^{m, p}_0(\Om)$, which contains functions for which all partial
derivatives of order lower or equal to  $m$ belong to $L^p$ (the
subscript 0 means null boundary conditions). Whithin these spaces, the
easier to use (both in theory and practice) are the Hilbert spaces
$W_0^{m,2}$, that can be equipped with the norm
$$
\|v\|_m^2 = \sum_{|\al| = m} \int_\Om
\left(\frac{d^mv}{dx_1^{\al_1}\ldots dx_n^{\al_n}}\right)^2 dx
$$
the sum being over all multi-indices $(\al_1, \ldots,
\al_n)$ such that $|\al| = \al_1 + \cdots+ \al_n = m$. The smallest
$m$ that makes this norm admissible depends on the dimension. One can
take $m=2$ in dimension 1, $m=3$ in dimension 2
and 3.

In the last case, the space $L^2([0,1], \CB)$ also is a Hilbert space,
with norm
$$
\|\bfv\|_{2, \CB}^2 = \int_0^1 \|\bfv(t, .)\|^2 dt
$$
and
$$
d_\CB(g, g') = \inf\{ \|\bfv(., .)\|_{2, \CB}: \bfg_\bfv(1,g(x)) \equiv g'(x)\}
$$

For Hilbert spaces, bounded sets are weakly relatively compact, and
the norm is lower semi-continuous. This implies that, if 
$\bfv_n$ is a mminimizing sequence for the distance  $d_\CB$, that is
if  $\|\bfv_n\|\ria
d_\CB(g, g')$ and for all $n$, $\bfg_{\bfv_n}(1,g(x)) \equiv
g'(x)$, then there exists a sub-sequence of  $\bfv_n$ that weakly converges
in $L^2([0,1], \CB)$ to some $\bfv^*\in\CB$ such that
$\|\bfv^*\| \leq d_\CB(g, g')$. One can prove that weak convergence
for the Hilbert structure is a stronger condition than the one given
in definition
\ref{weak*}; from theorem \ref{weak.closed}, we get the fact $\bfg_{\bfv^*}(1,g(x)) \equiv
g'(x)$. This implies that $d_\CB(g, g') = \|\bfv^*\|$: $\bfv^*$
provides an optimal matching trajectory between $g$ and $g'$. The
existence of such an optimal trajectory is very important in practice,
since, as we will see later, it is closely related to the issue of
aligning deformable objects.


\chapter{Estimation of Diffeomorphisms}
\label{est.diff}
\section{Introduction}
In the previous chapters, we have already noticed the importance of
performing a correct registration betwee deformable objects.
Such a registration can formally be interpreted as estimating a
diffeomorphism between the compared objects.

This chapter provides several algorithms that achieve this task. It is
organized as follows. We first provide a list of measures that
evaluate the ``size'' of a diffeomorphism. These definitions will
provide essential components for matching techniques, which will be
presented in the second part of the chapter.


\section{Quantitative Evaluation of Diffeomorphisms}
\label{size}
\subsection{Definitions}
We first introduce some notation. We let $\Om$ be a bounded open set
in $\mR^d$. We will consider a group $G$ of
diffeomorphisms (like one of those proposed in the previous chapter)

This group acts on {\em deformable structures} that will correspond to
shapes or other visiaull objects. For example, one can define the
action on configurations of points and on images as follows

\medskip
\encadre{
\begin{definition}
\label{action}
If $\bfx = (x_1, \ldots, x_N) \in
\Om^N$, and $g\in G$, we let
$$
g.\bfx = (g(x_1), \ldots, g(x_N))
$$
If $I$ is an image defined on $\Om$ and $g\in G$, we define $g.I$ by
$$
g.I(x) = I \circ g^{-1}(x)
$$
\end{definition}
}
\medskip

It is easy to check that these define actions. Let us do it for
images, since the case of landmarks is really straightforward. We want
to prove $g.(h.I) = (g\circ h).I$. But the former is $g.(I\circ
h^{-1}) = I \circ h^{-1} \circ g^{-1} = I \circ (g\circ h)^{-1} =
(g\circ h).I$. We will study these actions (and some other) in more
detail in the next chapter. For now, our goal is to quantify the size
of an element $g$ in $G$, by a quantity denoted $E(g)$, which will
measure, in some way, how $g$ differs from the identity mapping. The
idea, in order to quantify the variation of deformable
objects, will be to use this measure to evaluate the amount of
distorsion that is required to align them.

\subsection{Standard norms on functions}
One approach, to evaluate the signe of $g\in G$, is to define $u = g-
\id$ and use one of the numerous norms provided by functional analysis
to measure $u$. We cannot provide a complete outlook of these norms
here, but present some important families of examples.

We first start with norms of Sobolev type, that is, involving the
$L^p$ norms of some derivatives of $u$. The simplest example, say in
dimension 1, is to let
$$
E(g)=\int_\Om (g_x-1)^2 dx\,
$$
with some suitable boundary condition (e.g. $u=0$ at the extrimities
of $\Om$.

More generally, we can take, this time in any dimension
$$
E(g) = \|u\|_L^2 = \scp{Lu}{u}_2 = \int_\Om \scp{Lu(x)}{u(x)} dx\,.
$$
where $L$ is an operator defined on a space $\CH_L\subset \CL^2(\Om,
\mR^n)$, with the assumption that $\scp{Lu}{u} \geq c \norm{u}_2^2$,
for some $c>0$
($L$ said to be strictly monotonous). One of the advatages of this
last assumption is that it allows one to complete $\CH_L$, equipped
with $\|.\|_L$, into a
Hilbert space embedded in $L^2$ (this is the Friedrich extension,
cf. \cite{zei95}). One simple choice is to take $L = (\al - \Delta)^k$
for some $k>0$. For such norms, Sobolev inclusion theorem are
typically used to control the smoothness properties of elements of $\CH_L$.


\bigskip
One can also build a norm using an orthonormal basis of  $\CL^2(\Om,
\mR^k)$, that we shall denote $(f_1, \ldots, f_n, \ldots)$. Any $u$ in
$L^2(\Om, \mR^k)$ can be decomposed under the form
$$
u = \sum_k u_k f_k
$$
with
$$
u_k = \int_\Om u(x) f_k(x) dx.
$$
Letting $c = (c_1, \ldots, c_k, \ldots)$ be a sequence of positive numbers,
we can set
$$
\|u\|_c^2 = \sum_{k=1}^\infty c_k u_k^2\,.
$$

When the $c_k$'s are bounded, this will converge for all $u\in
\CL^2(\Om, \mR^k)$. But this is not the interesting situation; when
the $c_k$'s tend to infinity, the set $H_c$ that contains all $u$'s such
that $\|u\|_c$ is bounded is a strict subspace of $L^2(\Om, \mR^k)$,
which forms a Hilbert space when equipped with $\|.\|_c$. In this case,
for an image $u$ to belong to $H_c$, its high order coefficient must
tend to 0 faster than they would have had under the usual $L^2$
constraint which requires that the sum of squares of the $u_k$
converge. Because orthonormal bases of $L^2(\Om, \mR^k)$ are usually
ordered so that high order coefficients analyze high frequencies of
the function, a function in $H_c$ will generally be a smooth
function. For several bases (like Fourier, wavelets), the smoothness
of $u$ can be deduced from the rate of convergence of its coefficients
to 0   \cite{mey92}. 


There is another advantage in using such a decomposition, which comes
from the fact that it
provides a graded representation of the function. By truncating this
decomposition, one can obtain a good approximation of $u$. When
solving an optimization problem involving $u$, it becomes easy to
start clamping high order coefficients to zero, and progressively
release them during the procedure. This provides a coarse-to-fine
algorithm that is often computationally efficient, and also
potentially avoids being trapped in early local minima. One of the
problems with this representation, however, is that the evaluation of
$u$ at a point $x\in \Om$ can sometime be expensive, since it requires
summing the series that decomposes $u$. Fast transform exits, however,
for this purpose, on the most popular orthonormal bases (Fourier, wavelets).


\bigskip
These two points of view are equivalent when the basis $(f_k)$
diagonalizes the operator $L$. In such a case, letting
$L f_k = c_k f_k$, we have
$$
\int_\Om |Lu|^2 = \sum_{k=1}^\infty c_k u_k^2.
$$
Such a decomposition has been used to optimize deformations, for
example in
 \cite{ap91, dss96, ass98}, \cite{tro98}.


\subsection{Hyperelastic models}

\index{hyperelasticity}
The theory of elasticity uses a few ``universal'' principles,
like material indifference, or isotropy to derive general constraints
on the evaluation of the deformation that is applied on an object. To
these principles, that remain valid for many shape recognition
problems, one adds the constraint the the elastic energy only depens
on the first derivative of the deformation. Such constraints naturally
lead to the definition of a hyperelastic material; if   $g$ is a
diffeomorphism defined on  $\Om$, its associated deformation energy
must take the form 
$$
E(g) = \int_\Om W(x, C) dx
$$
where $C = {}^tDg.Dg$ is called the Cauchy stress tensor. For a
homogeneous material, $W$ does not depend on $x$. If the material is
also isotropic, $W$ only depends on the eigenvalues of $C$ (the
singular values of $Dg$). In this case, designing $W$ is simply
designing a function of $d$ variables (where $d$ is the dimension).

There are several classical definitions of such models. The simplest
nonlinear model is for  Saint-Venant Kirchhoff materials, for which (leting $E =
(C- I)/2$)
$$
W(C) = \frac{\la}{2} \left[\trace(E) \right]^2 + \mu \trace
\left(E^2\right).
$$
Here, $\la$ and $\mu$ are called the {\em Lam\'e} coefficients of the material (cf.
\cite{rwcm95}).

Most often, one uses the linearized version of this energy. If $g =
\id + u$ and $u$ is small, the approximation  $E
\simeq (du + {}^t du)/2$ can be made. Replacing $u$ by this
approximation yields an energy which is quadratic in $u$ and linear
Euler equations. An interesting matching algorithm based on this model
is presented in  \cite{ric02}.

One of the issues with Saint-Venant Kirchhoff materials is that their
energy does not penalize small eigenvalues of the tensor. This may
lead to singularities, and non-diffeomorphic solutions in large
deformation matching problems. To avoid this, one needs to penalize
small determinants of $C$. For example  \cite{hgygm99}, one can use
(in 2D) 
\begin{equation}
\label{elas.ax}
W(C) = W(\al, \be) = (\al^2+\be^2) (1+(\al\be)^{-2})
\end{equation}
$\al$ and $\be$ being the eigenvalues of $C$.

\subsection{Geodesic energies}
\label{viscous}
In chapter \ref{grp.diff}, we have considered diffeomorphisms $h$ for
which there exist solutions $\bfv$  of the boundary value problem
\begin{equation}
\label{ode}
\der{\bfg}{t}(t, .) = \bfv(t, g(t, .))
\end{equation}
with conditions $\bfg(0, .) = \id$ and $\bfg(1, .) = h$,
where $\bfv$ is a time dependent vector field such that
$$
\Ga(\bfv) := \int_0^1 \|\bfv(t,.)\|^2 dt < \infty,
$$
with the admissibility condition on $\|\bfv(t,.)\|$, as described in
section \ref{ext.diff}, chapter \ref{grp.diff}.

When they exist, the solutions $\bfv$ associated to a given $h$ are
not unique. It is natural to let the size of $h$ be the minimum of
$\Ga(\bfv)$ over all such $\bfv$; more precisely, letting $\CV_h$ be
the set of all time-dependent vector fields  $t\mapsto \bfv(t, .)$ on
$\Om$ such that the solution of the ordinary differential equation
$$
\der{y}{t} = \bfv(t, y)
$$
with initial condition $y(0) = x$ satisfies, for all $x$,  $y(1) =
h(x)$. We define
$$
E(h) = \inf_{\bfv\in \CV_h}\int_0^1 \|\bfv(t)\|^2 dt
$$

This may seem a rather awkward way to define the size of $h$. In
particular, defining it as an optimization problem may be
computationally inefficient. Also, for a given $h$, it is generally
difficult to decide whether  $\CV_h\neq \emp$, and even harder to
describe its elements.

This point of view, however, has several advantages. First, the
description of $\CV_h$ can be avoided by parameterizing the problem
dircetly in terms of  $\bfv$: instead of looking for the smallest $h$
achieving a given matching problem, one would look for the smallest
$\bfv$ such that the solution of the ODE at time  1 solves the
same matching problem. les m�mes contraintes. There is a numerical
cost, however, since this add a new dimension (time) to the
problem. But the additional complexity if balanced by the simplicity
of the cost function expressed in terms of $\bfv$. The continuous time
deformation
process (flow) that is obtained by integrating $\bfv$ also has some
interest, since it can provide smooth continuous evolution of one
initial deformable object to another
(cf. \cite{lv98}).

Anther advantage (that we will check later) is that the minimized cost
function can be interpreted as a metric distance (between the identity
and the function  $h$) on the group $G$. This diatance can in turn
induce distances between deformable objects. Symmetry and especially
triangular inequality are much harder to obtain with elastic
formulations.


\section{Landmark matching interpolation}

We start our review of specific matching problems with what is in fcat
an interpolation issue. In this setting, a certain number of point
correspondences are given, and the goal is to extrapolate them to a
global diffeomorphism on $\Om$. Unsurprisingly, the relevant methods
in this framework will borrow largely from the theory of interpolating
splines.

As before, $\Om$ is the support of the  configurations. We assume that
two sets of landmarks $(x_1, \ldots,
x_N)$ and $(y_1,
\ldots, y_N)$, bothe belonging to $\Om^N$. Our goal is to find a
diffeomorphism $g$ on $\Om$, with minimal size (as defined in the
previous section), such that, for all  $i$, $g(x_i) = y_i$ ({\em exact
  interpolation}), or,  $g(x_i) \simeq y_i$ ({\em inexact
  interpolation}). The last situation will be implemented using
penalities on the difference within an optimization problem.
Before providing a solution of these problems, we need to present some
required results for spline theory.
\index{splines}

\subsection{Splines} 

From an abstract point of view, spline approximation can be seen as
the application of the following very simple algebraic problem.
Let $\CH$ be a Hilbert space, and assume that $f_1, \ldots, f_N\in
\CH$, and $c_1, \ldots, c_N\in \mR$ are given. We introduce two
variational problems:
\begin{itemize}
\item[1.] Minimize $\|h\|$ over $h\in \CH$, under the constraints that
$\scp{f_i}{h}  = c_i$ for $i=1, \ldots, N$.
\item[2.] For a parameter $\la > 0$,  minimize $\|h\|^2 + \la \sum_{i=1}^N
(\scp{f_i}{h} - c_i)^2$
\end{itemize}
The two problems will respectively correspond to exact and inexact
interpolation.
The constraints, case 1., and the penalty in case 2., are not affected
when $h$ is replaced by $h+v$ and $v$ is othogonal to the finite
dimensional vector space generated by $f_1, \ldots, f_N$.

Introduce the $N$ by $N$ matrix $S$ with $S_{ij} =
\scp{f_i}{f_j}$. Acording to the discussion above, we can look for $h$ under the form
$$
h =\sum_{i=1}^N \al_i f_i,$$
for which $\|h\|^2 = \trp{\al} S \al$. 

Consider problem 1. We need to minimize $\trp{\al} S \al$ under the
constraints  $S\al = c$ ($\al$ and $c$ being vectors with coordinates
$\al_i$ and $c_i$). Assume that $S$ is invertible: then, the only
solution satisfying the constraints is $\al = S^{-1} c$ and there is
no minimization issue. When $S$ is not invertible, we must assume that
$c$ belongs to the renge of $S$, otherwise the problem will have no
solution. This implies that $c = S \al_0$ for some $\al_0$, and any
solution $\al$ must take the form $\al= \al_0 + \be$ with $\be \in
\text{ker}(S)$. But, for any such solution, we have $\trp{\al} S \al
= \trp{\al_0} S \al_0$ so that any $\al$ is a valid minimizer.

The fact that $S$ is not invertible is equivalent to the existence of
a non-vanishing $\be$ such that $\trp{\be} S \be = 0$, which can be
rewritten as $\|\sum_{i=1}^N \be_i f_i\|^2 = 0$. In other terms, $S$
is not invertible if and only if $f_1, \ldots, f_N$ are linearly
dependent. The solution $\al$ is this case is thus defined up
to
$\be$'s for which $\sum_{i=1}^N \be_i f_i=0$  which is  
obvious since adding such a linear combination does not change the
value of $h$.

Consider now problem 2. We need to minimize
$$\trp{\al} S \al + \la \trp{(S\al - c)}(S\al - c)$$
with respect to $\al$. This is a convex function of $\al$ and its
gradient is $S((\la S + I)\al - \la c)$. If $S$ is invertible, the
unique solution is $\al = S_\la c$ with $S_\la = S + I/\la$. In the
general case, it takes the form $S_\la^{-1}(c + \be)$ with $S\be = 0$.
Note that in this case, a solution always exists.
\bigskip

We now make the connection with spline interpolation. The goal here is
to find a smooth real-valued function $h$ � valeurs r�elle, defined on
$\Om$, such that  $h(x_i) = c_i$ (or
$h(x_i)\simeq c_i$) for some $x_i\in\Om$ and
$c_i\in\mR$. The smoothness of $h$ is measured by a Hilbert space
norm, for example
$$
\|h\|^2_L = \int_\Om \scp{Lh}{h}^2 dx
$$
where is is a strictly monotonous operator. We let $\CH_L$ be the
corresponding Hilbert space.

We can educe the linear constraints  $h(x_i) = c_i$ to those in the
abstract problem provided there exists, for each  $x_i$, an element
$f_{x_i}$ in $\CH_L$ such that, for all $h\in \CH_L$
$$
h(x_i) = \scp{f_{x_i}}{h}_L.
$$
By the Riesz representation theorem, this is equivalent to the fact
that the evaluation functional $h \mapsto h(x)$ is continuous for the
Hilbert structure. Spaces that satisfy this property for every $x$ are
called {\em reproducing kernel Hilbert spaces} (RKHS). The fact that a
space $\CH_L$ as above is a RKHS can typically be proed by invoking a
Sobolev inclusion theorem. We will henceforth assume that this is the
case, and therefore obtain the existence, for each $x$, of a function
$f_x\in \CH_L$ with, for all $h\in\CH_L$, $
h(x) = \scp{f_{x}}{h}_L.$ 

It can be noted that, taking $h=f_y$ in the previous formula yields 
$f_y(x) = \scp{f_{x}}{f_y}_L = f_x(y).$ The function of two variables
$(x,y) \mapsto K(x,y) = f_x(y)$ is the kernel of the space $\CH_L$,
which therefore satisfies the eproducing identity
$$
K(x,y) = \scp{K(x,.)}{K(y,.)}_L.
$$

Returning to our interpolation problem, we see that the constraints
now are  $\scp{f_{x_i}}{h}_L = c_i$ with $f_{x_i} = K(x_i, .)$ which
places us directly in the abstract situation we have studied. We
therefore obtain solutions $h$ of the form
$$
h(.) = \sum_{i=1}^N \al_i K(., x_i).
$$
The matrix $S$ being formed with the dot products of the $f_{x_i}$'s
is there the matrix with coefficients $S_{ij} = K(x_i, x_j)$, which
yields explicit formulas to compute $\al$ when $K$ is known.

\medskip

The availability of an explicit expression for $K$ is therefore an
essential practical requirement. When $L$ is a differential operator, there are very few situations in
which this kernel (which is also called the Green kernel of the
operator $L$) is computable. They generally are limited to cases for
which $L$ is some polynomial of the Laplacian, and to specific
boundary conditions on $\Om$ (or at infinity if $\Om$ is not bounded).
An important case of affine conditions at infinity provides thin-plate
splines, and will be presented in section \ref{bookstein}. For
operators of the form $(a - \Delta)^n$, the kernel can be computed via
the evaluation of Bessel functions. Explicit expressions for $L = -
\nabla ^n$ on rectangular domains with zero
boundary conditions is also available.

I one does not restrict to differential operators, however, there is
an explicit way to proceed, which consists in directly defining the
kernel $K$, provided one can prove that a valid $L$ exists without
necessarily having to compute it. Note that for the practical use of
the spline problem,
only $K$ is needed, not $L$.

Suficient conditions for $K$ can be taken as follows.
\begin{itemize}
\item Symmetry: $f_x(y) = f_y(x)$
\item $K$ is continuous and, for all $x\in\Om $
$$
\int_\Om K(x,y)^2 dy < \infty
$$
\item $K$ induces a positie operator on  $\CL^2$: for all $u\in \CL^2$
$$\int_\Om u(x) u(y) f_x(y) dx dy \geq 0.$$
\end{itemize}
One simple way to ensure the last condition is to take $K(x,y) =
F(x-y)$ where 
$F$ is the Fourier transform of an even, positive function; $K$ is, in
this case, a convolution kernel.

One can also restrict to radial kernels, which take the form $K(x,y) =
G(|x-y|)$ (\cite{ar95}, \cite{adry94}). For such a kernel to be valid
in every dimension, one can show that  $q(t)
= G(\sqrt{t})$
must be the Laplace transform of a meaure $\mu$ on $\mR^+$:
$$
q(t) = \int_0^{+\infty} e^{-\la u} d\mu(\la)\,.
$$
When $\mu$ is a Dirac measure, we obtain the important example of a
Gaussian kernel given by
$$
G(t) = \exp\left(-\frac{t^2}{\sig^2}\right).
$$
For more details, we refer to \cite{dyn89} and references therein.
\index{Gaussian kernel}



\subsection{Diffeomorphic interpolation}

\label{par:splines}
\subsubsection{Generalities}

We now extend spline theory to our diffeomorphic interpolation: extend
to $\Om$ a correspondence $x_i \to y_i, i=1, \ldots, N$. One important
remark, in this context, is that the interpolation to e computed is a
function from $\Om$ to $\Om$ and not from $\Om$ to $\mR$ as before. We
will therefore have to deal with several components. 


Before addressing the corespondence problem, we first start
with a direct generalization of the prevous section, in which we
consider a Hilbert space $\CH$ of smooth vector fields on $\Om$ and the
issue is to find a function $h$ in $\CH$ subject to
the constraints $h(x_i)= v_i$ for $x_1, \ldots, x_N\in\Om$ and $v_1,
\ldots, v_N \in\mR^d$. We will assume as before that the evaluation
function $h\mapsto h(x)$ is continuous.

The difference with the previous problem therefore is that the
function $h$ takes vector values, so that the function $h\to h(x)$ is
not a linear form anymore. Clearly, one can consider coordinate
functions and apply the previous theory. More generally, one can
associate to any vector $a\in \mR^d$ and any $x\in\Om$ the linear form
$\de_x^a: h \mapsto a^T h(x)$. By Riesz representation theorem, there
exists a vector field $f_x^a\in \CH$ such that $\scp{f_x^a}{h} = a^T
h(x)$ for all $h\in H$. Since $\de_x^a$ is linear in
$a$, the same property is true for $f_x^a$. Fixing $x,y\in \Om$, we
therefore obtain the fact that $a \mapsto f_x^a(y)$ in linear from
$\mR^d$ to $\mR^d$. It can therefore be put under the form
$f_x^a(y) = K(y, x)a$ where $K(y, x)$ is a $d$ by $d$ matrix. The
matrix-valued function $K$, defined on $\Om^2$, if the reproducing
kernel of the space $\CH$. It is characterized by the property that,
for all $h\in \CH$, for all $x\in\Om$ and $a \in \mR^d$:
$$
a^Th(x) = \scp{K(x,.)a}{h}_V.
$$
This implies in particular that, for all $x$ and $a$, $K(., x) a \in
\CH$. Moreover, we can write, for all $x,y\in\Om$, for all
$a,b\in\mR^d$:
$$
\scp{K(., x)a}{K(., y)b}_\CH = a^T K(x,y) b
$$
which is the $d$-dimensional reproducing property of the kernel
$K$. This implies in particular the symmetry property: $K(x,y) =
K(y,x)^T$ for all $x,y\in\Om$.

Returning to our interpolation problem, each constraint $h(x_i) = v_i$
is equivalent to $e_k^T h(x_i) = v_{ik}$ for $k=1, \ldots, d$, where
$e_1, \ldots, e_d$ is the canonical basis of $\mR^d$ and $v_i =
(v_{i1}, \ldots, v_{id})$ is the decomposition of $v_i$ in this
basis. Since $e_k^T h(x_i) = \scp{K(., x_i)e_k}{h}_\CH$, we are reduced
to our previous abstract problem, this time with $d N$ constraints
instead of $N$. We obtain, in particular the fact that an optimal
solution for the exact and inexact interpolation problems can be
written under the form, for some real numbers $\al_{ik}$,
$$
h(.) = \sum_{i=1}^N \sum_{k=1}^d \al_{ik} K(., x_i) e_k
$$
which can also be expressed as (letting $al_i = (\al_{i1}, \ldots, \al_{id})^T\in\mR^d$)
$$
h(.) = \sum_{i=1}^N K(., x_i) \al_i.
$$

The solution of the exact matching problem now can be obtained by
solving the linear system
$h(x_j) = \sum_{i=1}^N K(x_j, x_i) \al_i$. Introducing the $dN$ by
$dN$ matrix $S$, organized in $d$ by $d$ blocs, with $i, j$ bloc given
by $K(x_i, x_j)$, and assuming that $S$ is invertible, the solution if
the the system is
$$
\begin{pmatrix} \al_1\\ \vdots \\ \al_N\end{pmatrix} = S^{-1}
\begin{pmatrix} h(x_1)\\ \vdots \\ h(x_N)\end{pmatrix}.
$$

The solution of the inexact matching problem, namely minimize
$$
\|h\|_{\CH} + \la \sum_{i=1}^N |h(x_i) -  v_i|^2
$$
is given by the same expression with $S$ replaced by $S_\la  = S +
\Id/\la$. 

The systems above becomes simpler uneder the assumption, which is
widely made in practice, that the kernel is diagonal, i.e.,
$K(x,y) = k(x,y) \Id$ with $k(x,y)\in\mR$. In this case, the previous
solution can be reordered as
$$
\begin{pmatrix} \al_{11} \ldots \al_{1d} \\ \vdots \\ \al_{N1} \ldots
\al_{Nd} \end{pmatrix} = \tilde S^{-1}
\begin{pmatrix} h_1(x_1)\ldots h_d(x_1) \\ \vdots \\ h_1(x_N)\ldots
h_d(x_N) \end{pmatrix}
$$
where $\tilde S_{ij}  = k(x_i,x_k)$ is a $N$ by $N$ matrix. This
reduces the issue of solving a $dN$-dimensional system to solving $d$
$N$-dimensional systems, which is much less expensive.

\subsubsection{Invariance of the inner-product}

We haven't discussed so far under which criteria the dot product (or
the operator $L$) should be selected. One important criterion, in
particular for shape recognition is equivariance or invariance with
respect to rigid transformations. We here consider the situation $\Om
= \mR^d$.

Our first requirement is as follows. Let $x_1, \ldots, x_N$ be the
landmarks, and $v_1, \ldots, v_N$ the vectors. We want the obtained
interpolation $h$ to be invariant by an orthogonal change of
coordinates. This means that, if $A$ is a rotation and $b\in\mR^d$, and if we consider the problem associated
to the landmarks $Ax_1+b, \ldots. Ax_N+b$ and the constraints $Av_1,
\ldots, Av_N$, the obtained function $h^{A,b}$ must satisfy
$h^{A,b}(Ax+b) = A h(x)$ or $h^{A,b}(x) = A h (A^{-1}x - A^{-1} b) := ((A,b)\star h)(x) $.
The transformation $h \to (A,b)\star h$ is an action of the group of
translations-rotations on functions defined on $\Om$. To ensure that
$h^{A,b}$ always satisfies this condition, it is sufficient to enforce
that the norm is invariant by this transformation: $\|(A,b)\star h\|_\CH  = \|h\|_\CH$.

So, we need the action of translation and rotations to be an isometry
in $\CH$. Assume that $\|h\|_\CH^2 = \lform{Lh}{h}$ for some operator
$L$, and define the operator $A*L$ the fact that, for all $h, \tilde h\in
\CH$ 
$$
\lform{L((A,b)\star h)}{(A,b)\star \tilde h} = \lform{((A,b)\star L) h}{\tilde h}
$$
so that, the invariance condition becomes expressed in terms of $L$:
$(A,b)\star L  = L$. 

Denoting $h_k$ the $k$th coordinate of $h$, and
$(Lh)_k$ the $k$th coordinate of $Lh$, we can write the matrix
operator $L$ under the form
$$
(Lh)_k  = \sum_{l=1}^d L_{kl} h_l
$$
where $L_{kl}$ is a scalar operator.  We first consider the case when
eack $L_{kl}$ is a pseudo-differential operator in the sense that, for
any $C^\infty$ scalar function $u$, the Fourier transform of $L_{kl} u$ takes the
form, for $\xi\in\mR^d$:
$$
\hat{L_{kl} u} (\xi) = f_{kl}(\xi) \hat u(\xi) 
$$
for some scalar function $f_{kl}$. An interesting case is when
$f_{kl}$ is a polynomial, which corresponds, from standard properties
of the Fourier transform, to $L_{kl}$ being a differential
operator. Also, by the isometric property of Fourier transforms,
\begin{eqnarray*}
\lform{Lh}{h} &=& \sum_{k=1}^d \lform{(Lh)_k}{h_k} \\
&=& \sum_{k, l=1}^d \lform{L_{kl} h_l}{h_k} \\
&=& \sum_{k, l=1}^d \lform{\hat{(L_{kl} h_l)}}{\hat h_k} \\
&=& \sum_{k, l=1}^d \lform{f_{kl} \hat h_l}{\hat h_k}.
\end{eqnarray*}
Note that since we want $L$ to be a positive operator, we need the
matrix $f(\xi) = (f_{kl}(\xi))$ to be a complex, positive, Hermitian
matrix. Moreover, since we want $\lform{Lh}{h} \geq c \lform{h}{h}$
for some $c>0$, it is natural to enforce that the smallest eigenvalue
of $f(\xi)$ is uniformly bounded from below (by $c$). 

To use this expression, we need to compute $\hat{(A,b)\star h}$ in
function of $\hat h$. We have
\begin{eqnarray*}
\hat{(A,b)\star h}(\xi) &=& \int_{\mR^d} e^{-i\xi^T x} (A,b)\star h(x)
dx\\
&=& \int_{\mR^d} e^{-i\xi^T x} Ah(A^{-1} x - A^{-1} b)
dx\\
&=& \int_{\mR^d} e^{-i\xi^T (Ay+b) } Ah(y) dy\\
&=& e^{-i\xi^T b}A h(A^T \xi)
\end{eqnarray*}
where we have used the fact that $\det A = 1$ (since $A$ is a rotation
matrix) in the change of variable. We can therefore write, denoting
$\phi_b$ the function $\xi \mapsto \exp(-i\xi^T b)$
\begin{eqnarray*}
\lform{L(A\star h)}{A\star h} &=& \sum_{k, l=1}^d \sum_{k', l'=1}^d
\lform{\phi_b f_{kl} a_{ll'} \hat h_{l'} \circ A^T}{\phi_b a_{kk'} \hat
h_{k'}\circ A^T}\\
&=& \sum_{k, l=1}^d \sum_{k', l'=1}^d
\lform{ f_{kl} \circ A  a_{ll'} \hat h_{l'}}{a_{kk'} \hat
h_{k'}}.
\end{eqnarray*}
We have used the fact that $\det A = 1$ and that $A^T = A^{-1}$. The
function $\phi_b$ cancels in the complex product:
$$
\lform{u}{v} = \int_{\mR^d} u \bar v d\xi
$$
where $\bar v$ is the complex conjugate of $v$.

Reodering the last equation, we obtain 
$$
\lform{L(A\star h)}{A\star h} = \sum_{k', l'=1}^d
\lform{ \sum_{k, l=1}^d a_{kk'} f_{kl} \circ A  a_{ll'} \hat h_{l'}}{ \hat
h_{k'}}
$$
which yields our invariance condition in terms of the $f_{kl}$: for
all $k,l$
$$
f_{kl} = \sum_{k', l'=1}^d a_{k'k} f_{k'l'} \circ A  a_{l'l}
$$
or representing $f$ as a hermitian matrix $A^T f(A\xi) A = f(\xi)$ for all $\xi$ and
all rotation matrix $A$. 

Consider first the case $\xi=0$. The identity $A^Tf(0) A = f(0)$ for
all $A$ implies, by a standard linear algebra theorem, that $f(0) =
\la \Id$ for some $\la \in \mR$.  
We now fix $\xi\neq 0$, and let $\eta_1 = \xi/|\xi|$. Extend $\eta_1$
into an orthonormal basis, $(\eta_1, \ldots, \eta_d)$ and let $(e_1,
\ldots, e_d)$ be the canonical basis of $\mR^d$. Define $A_0$ by $A_0
\eta_i = e_i$ so that $A_0\xi = |\xi| e_1$ and
$f(\xi) = A_0^T f(|\xi| e_1) A_0 = M_0$. Now, for any  rotation $A$ that
leaves $eta_1$ invariant, we also have
$f(\xi) = A^T A_0^T f(|\xi| e_1) A_0 A$ since $A_0 A\eta_1 =
e_1$. This yield the fact that, for any $A$ such that $A \eta_1  =
\eta_1$, we have $A^T M_0 A = M_0$.

This property implies, in particular that $A^T M_0\eta_1 = M_0 \eta_1$
for all $A$ such that $A\eta_1 = \eta_1$. Taking a decomposition of
$M_0\eta_1$ on $\eta_1$ and $\eta_1^\perp$, it is easy to see that
this implies that $M_0 \eta_1 = \al \eta_1$ for some $\al$. So
$\eta_1$ is an eigenvector of $M_0$, which implies that $\eta^\perp$
is closed by the action of $M_0$. But the fact that $A^T M_0 A  = M_0$
for any rotation $A$ of $\eta_1^\perp$ implies that $M_0$ restricted
to $\eta_1\perp$ is a homothetie, i.e., that there exists $\la$ such
that $M_0\eta_i = \la \eta_1$ for
$i\geq 2$. 

Using the decomposition $u= \scp{u}{\eta_1} \eta_1 + u -
\scp{u}{\eta_1} \eta_1$, we can write
$$
M_0u= \al \scp{u}{\eta_1} \eta_1 + \la (u -
\scp{u}{\eta_1} \eta_1) = ((\al-\la) \eta_1 \eta^T  + \la \Id) u
$$
which finally yields the fact that (letting $\mu = \al- \la$ and using
the fact that the eigenvalues of $M_0 = A^T f_0(\|xi\| e_1) A$ are the
same as those of $f_0(|\xi|e_1)$ and therefore only depend on $|\xi|$) 
$$
f(\xi) = \mu(|\xi|) \frac{\xi\xi^T}{|\xi|^2} + \la(|\xi|)\Id.
$$
Note that, since we know that $f(0) = \la\Id$, we must have $\mu(0) =
0$.

Let's first consider the case $\mu = 0$: in this case, we have $f_{kl}
= 0 $ if $k\neq l$ and $f_{kk} (\xi) = \la(|\xi|)$ for all $k$, which
implies that $L_{kl} = 0$ for $k\neq l$ and all $L_{kk}$ are equal to
a fixed scalar operator $L_0$, i.e.,
$$
Lh = (L_0h_1, \ldots, L_0 h_d).
$$

This is the type of operators which is mostly used in practice. If one
wants $L_0$ to be a differential operator, $\la(|\xi|)$ must be a
polynomial in the coefficients of $\xi$, which is only possible  if it
tales the form
$$
\la(|\xi|) = \sum_{i=0}^p \la_p |\xi\|^{2p}
$$
in which case the corresponding operator is $L_0 u = \sum_{i=0}^p
(-1)^p \Delta^p u$. For example, if $\la(|\xi|) = (\al + 
|\xi|^2)^2$, we get $L_0 = (\al I - \Delta)^2$. 

Let's consider an example of valid differential operator with $\mu\neq
0$, and take $\mu(|\xi|) = \al |\xi\|^2$ and $\la$ as above. This
yields the operator 
$$L_{kl} u  =  - \al   \frac{\prt^2 u}{\prt x_k\prt x_l} + \de_{kl} \sum_{i=0}^p
(-1)^p \Delta^p u$$
so that
$$
Lh = - \al \nabla \text{div} h + (\sum_{i=0}^p
(-1)^p \Delta^p )h.
$$

\bigskip

The Kernel (or Green function) associated to an operator associated to
a given $f$ can be computed simply in the Fourier domain. It is indeed
given by
$$
\hat{K_{kl} u} (\xi) = g_{kl}(\xi) \hat u(\xi) 
$$
with $g= f^{-1}$ (the matrix inverse), as can be checked by a simple
computation. Note that, for the spline application, we need $K\de_x^a$
to be well defined. The $k$th component of $\de_x^a$ being $a_k
\de_x$, we can write $\hat {(\de_x^a)_k} (\xi) = a_k \exp(-i\xi^T x)$
so that
$$
\hat{(K\de_x^a)_k} = \sum_{l=1}^d g_{kl}(\xi) a_l \exp(-i\xi^T x).
$$
A sufficient condition for this to provide a well-defined function is
that all $g_{kl}$ are integrable , in which case
$$
(K\de_x^a(y))_k = \sum_{l=1}^d \tilde g_{kl}(y-x) a_l
$$ 
where $\tilde g$ is the inverse Fourier transform of $g$.
In the case $f = \la \Id + \mu \eta \eta ^T$, we have 
$$
f^{-1} = \frac{1}{\la} (\Id - \frac{\mu}{\la + \mu} \eta \eta^T).
$$

Let $\rho(\xi) = 1/\la(|\xi|)$ and $\psi(\xi) = \mu(\xi) / (\la(\xi)
|\xi|^2 (\la(\xi) +
\mu(\xi)))$, so that
$f^{-1}(\xi) = \rho(|\xi|) \Id - \psi(|\xi|) |\xi|^2 \eta
\eta^T$. The eigenvalues of $f^{-1}$ therefore are $\rho(\xi)$ and
$\rho(\xi) - \psi(\xi) |\xi|^2$. Since the eigenvalues of $f$ must be
bounded from below by $c$, we obtain the following constraints on
$\rho$ and $\psi$:
$$ |\xi|^2 \max(\psi, 0) \leq \rho \leq c. $$ 

We can moreover obtain the expression of the Kernel in terms of the
inverse Fourier transforms of $\rho$ and $\psi$, which is:
$$
K(x,y)_{kl} = \tilde g_{kl}(x-y) = \de_{kl} \tilde \rho (x-y) + \pddr{\tilde\psi}{x_k}{x_l}(x-y).
$$

The inverse Fourier transform of a radial function is a radial
function. We have in fact
\begin{eqnarray*}
\tilde\rho(x) &=& \int_{\mR^d} \rho(|\xi|) e^{i\xi^T x} d\xi \\
&=& \int_0^\infty \rho(t) t^{d-1} \int_{S^{d-1}} e^{i\eta^T t x}
ds(\eta) \\
&=& \int_0^\infty \rho(t) t^{d-1} \int_{S^{d-1}} e^{i\eta_1 t |x|}
ds(\eta) 
\end{eqnarray*}
where $S^{d-1}$ is the unit sphere in $\mR^d$, and the last identity
comes from the fact that the integral is invariant by roration of $x$,
which allows us to replace $x$ by $|x|$. The last integral can in trun
be expressed in terms of the Bessel function $J_{d/2}$ (\cite{so97}),
which will not be used here.

There therfore exists two functions $\ga_1$ and $\ga_2$ such that
$\rho(x) = \ga_1(|x|^2)$ and $\psi(x) = \ga_2(|x|^2)$ and the
expression for $K$ becomes
$$
K(x,y) = (\ga_1(|x-y|^2) + 2 \ga'_2(|x-y|^2)) I + 2\ga''_2(|x-y|^2) (x-y)(x-y)^T.
$$

Again, the choice $\ga_2 = 0$ is commonly made in practice, yielding
diagonal radial kernels \cite{dyn89}.
\index{radial kernels}



\subsection{Thin plate interpolation}
\label{bookstein}
\index{Thin plates}

Thin plate theory corresponds to the situation in which the operator
$L$ is some power of the Laplacian. As a tool for shape analysis, it
is been originally described by  Bookstein
\cite{boo91}. Consider the following bilinear form
\begin{equation}
\label{lapl}
\scp{f}{g}= \int_{\mR^2} \Delta f \Delta g dx\,.
\end{equation}

We need to define it on a Hilbert space which is somewhat
unusual. We consider the space $\CH_1$ of all functions in $\mR^d$ with
square integrable second derivatives, which have a bounded gradient at
infinity. In this space, $\|f\| = 0$ is equivalent to the fact that $f$
is affine, i.e., $f(x) = a^Tx + b$, for some $a\in \mR^d$ and
$b\in\mR$. The Hilbert space we consider is the space of $f\in \CH$
modulo the affine functions, i.e., the set of equivalent classes
$$
[f] = \defset{g: g(x) = f(x) + a^T x + b, a\in\mR^d, b\in\mR}.
$$ 
For $f\in\CH_1$. Obviously the norm associated to \eqref{lapl} is
constant over $[f]$ and $\|f\| = 0$ if and only if $[f] = [0]$ so that

This space also has a kernel, although the analysis has to be
different from what we have done previously, since the evaluation
functional is not defined on $\CH$ (functions are only known up to the
addition of an affine term). However, the following is true 
\cite{mei79}. Let $ U(r) = (1/8\pi) r^2\log r$ if the dimension, $d$, is 2, and
$U(r) = (1/16\pi) r^3$ for $d=3$. Then, for all $u\in\CH_1$, there
exists $a_u\in\mR^d$ and $b_u\in\mR$ such that 
$$
u(x) = \int_{\mR^2} U(|x-y|) \Delta^2 u(y) dy + a_u^Tx + b_u.
$$
We wil denote $U_x$ the function $y\mapsto U(|x-y|)$.

Fix, as in the previous section, landmarks $x_1, \ldots, x_N$ and
function values  $(c_1, \ldots,
c_N)$. Again, the constraint $h(x_i) = c_i$ has no meaning in $\CH$,
but the constraint  $\scp{U_{x_i}}{h}_\CH = c_i$ does, and means
$$
h(x_i) + a_h^Tx_i + b_h = c_i,
$$
i.e., $h$  satisfies the constraints un to an affine term.

Denote, as before
$S_{ij} =  U(|x_i-x_j|).$
The function $h$ which is optiml under the constraints
$\scp{U_{x_i}}{h}_\CH = c_i, i=1, \ldots N$ will therfore take the form
$$
h(x) = \sum_{i=1}^N\al_i U(|x-x_i|) + a^Tx + b.
$$
The corresponding interpolation problem is: minimize $\|h\|$ under the
constraint that there exists $a,b$ such that $h(x_i) + a_h^Tx_i + b_h
= c_i, i=1, \ldots, N$.The inexact matching problem simply consists in
minimizing
$$
\|h\|^2_\CH + \la \sum_{i=1}^N (h(x_i) + a^Tx_i + b - c_i)^2
$$
with respect to $h$ {\em and} $a,b$. 

Replacing $h$ by its expression in terms of the $\al$'s, $a$ and $b$
yields the finite dimensional problem: minimize
$$\al^T S \al$$
under the constraint  $S\al + Q\ga = c $ with $\ga = (a_1, \ldots, a_d,
b)^T$ (with size $d+1 \ti 1$) and $Q$, with size $N\ti d+1$ given by
(letting $x_i = (x_i^1, \ldots, x^d_i)$):
$$
Q = \left(\begin{matrix}
x_1^1& \cdots & x_1^d&1\\
\vdots&\vdots&\vdots\\
x_N^1&\cdots & x_N^d&1
    \end{matrix}\right).
$$
The optimal   $(\al, \ga)$ can be computed by identifying the gradient
to 0. One obtains
$$
\hat{\gamma} = \left(Q^T S^{-1}Q\right)^{-1} Q^T S^{-1}c
$$
et
$\hat{\al} = S^{-1}(c - Q\hat{\ga})$.

For the inexact matching problem, one minimizes
$$\al^T S \al + \la{}(S\al + {}^tax_i + b - c)^T(S\al + {}^tax_i + b
- c)\,,$$
and the solution is provided by the same formulae, simply replacing
$S$ by $S_\la = S + (1/\la)I$.

When the function $h$ (and the $c_i$) take $d$-dimensional values
(e.g., correspond to displacements), the above computation has to be
applied to each coordinate, which simply corresponds to using a
diagonal operator, each component equal to $\Delta^2$, in the
definition of the dot product.

The figure \ref{book.ex} provides an example of such an 
interpolation, togother with an interpolation with radial functions (\cite{ar95}).
\begin{figure}
\begin{center}
\includegraphics[width=4cm]{TUTORIALPS/one.ps}\ \
\includegraphics[width=4cm]{TUTORIALPS/two.ps}\\
\includegraphics[width=4cm]{TUTORIALPS/onetotwo.ps}\ \
\includegraphics[width=4cm]{TUTORIALPS/twotoone.ps}\\
\includegraphics[width=4cm]{TUTORIALPS/onetotwoG.ps}\ \
\includegraphics[width=4cm]{TUTORIALPS/twotooneG.ps}\\
\end{center}
\caption{\label{book.ex} Exact landmark interpolation. First Row:
landmaks placed on original images; Second Row: thin-plate
interpolation, form the first image to the second and vice-versa;
Third Row: interpolation using radial functions, with $K(x, y) = log(|x-y|^2 + c)$
(images take from the olivetti face database).}
\end{figure}

\medskip


When thin-plate spline methods (or the spline interpolation methods)
are applied to displacement field, there is no constraing that
guarantees that the obtained function, $\id + h$ is a
diffeomorphism. It is indeed not rare to observe folds in the
obtained solution. The following section describe a solution to
diffeomorphic interpolation.


\subsection{Diffeomorphic interpolation}
\label{joshi}

The distance we have defined on groups of diffeomorphisms (chapter
\ref{grp.diff}) can be used to interpolate landmark correspondences
\cite{jos97}. In this framework we introduce a time-dependent velocity
field on $\Om$, denoted   $t\mapsto v(t, .)$, and the ODE
$$
\der{y}{t} = \bfv(t, y)
$$
Let $\bfg_\bfv(t, x)$ be the solution of this equation at time $t$
with initial condition $\bfg_\bfv(0, x) = x$. The exact matching
problem consists in finding 
$\bfv$ to minimize
$$
E(\bfv) = \int_0^1 |\bfv(t)|^2_{\CH} dt
$$
under the constraints 
 $\bfg_\bfv(1, x_i) = y_i$ for $i=1, \ldots, N$.

Introduce the time dependent displacements, at each landmark, given by
$u_i(t) = \bfg_\bfv(t, x_i) $; these trajectories are unknown,
but once they are obtained, $\bfv$ can be reconstructed in a simple
way. Indeed, let's fix  $u_i(t), i=1, \ldots N, t\in[0, 1]$. First
notice that we have
$$
\der{}{t} u_i(t) = \bfv(t, u_i(t))
$$
so that specifying the $u_i$'s also specifies the values of 
$\bfv(t,  .)$ at $u_i(t)$ (it is in fact equivalent, when
combined with the initial condition $u_i(0) = x_i$). The remaining
problem therefore boils down to: minimize $E(\bfv)$ under the
constraints that $\bfv(t, u_i(t)) = \prt u_i /\prt t$. This
problem can be solved separately for each time, 
 i.e., minimize, for all $t$, $\|\bfv(t)\|_{\CH}$ under the
constraints. This is exactly the kind of problem that have been solved
in the previous sections. Let $(x,y) \mapsto K_x(y) =
K(x,y) $ be the reproducing kernel for $\CH$. Let $S(\bfu(t))$ be the bloc
matrix composed with $s_{ij}(t) = K(u_i(t),
x_j+u_j(t))$, we have, stacking the $u_i$ in a coloumn vertor $\bfu$,
$$
\bfv(t, x) = \sum_{i=1}^N  K_{u_i(t)}(x) \al_i(t)
$$
with $\al(t) = S(\bfu(t))^{-1} \der{\bfu}{t}$.
Again, when $K$ is diagonal, the computation can be split in $d$ $N$
dimensional systems.

Using this we obtained a residual energy which only depends on the
evolution of the $Nd$-dimensional vector $\bfu$:
 L'�nergie obtenue est
$$
E(\bfv) = E(\bfu) = \int_0^1 {}^t\der{\bfu}{t} S(\bfu(t))^{-1} \der{\bfu}{t} dt
$$
which needs to be minimized with constraints $u_i(0) = x_i$ and
$u_i(1) = y_i$. Several algorithms have been proposed for this purpose
\cite{jos97,
vmty04, aty05} and we refer the reader to these references for
implementation details. It is quite easy to come up with examples in
which the spline approach fails whereas the diffeomorphic
interpolation provides a smooth diffeomorphism
(Fig. \ref{fig:splines}).

\begin{figure}
\begin{center}
\includegraphics[width=0.4\textwidth]{FIGURES/init.eps}
\includegraphics[width=0.4\textwidth]{FIGURES/final0.eps} \\
\includegraphics[width=0.4\textwidth]{FIGURES/final.eps}
\includegraphics[width=0.4\textwidth]{FIGURES/trajectories.eps}
\caption{\label{fig:splines} Comparison between linear and
diffeomorphic landmark interpolation. First Row: Initial landmarks and
their displacements (Left); Deformation obtained by linear splines
(Right). Second Row: Deformation obtained by diffeomorphic inexact
matching, using the method proposed in \cite{aty05} (Left); Optimal
trajectories ($u_i(t)$) provided by diffeomorphic inexact landmark
matching.}
\end{center}
\end{figure}


\section{Matching Unlabelled Points}

\subsection{The assignment problem}

\subsection{Assignment combined with thin plates}

\subsection{Measure matching}

\section{Matching Images}

Matching images is somewhat different from matching sets of points that we
have considered so far. The matching procedure is not driven here by
the information of which point or set is matched to another, but by
the fact that some numerical quantity (a signature) is transported by
the deformation. This signature can be gray values for scalar images,
but can also, as we shall see, have other interpretation.

We will consider, here again, a bounded open subset $\Om \sub \mR^d$,
and images $I$ will be scalar or vector-valued functions $f:\Om\to
\mR^k$. If a deformation $g$ acts on $\Om$, each points $x$ is
transformed into $g(x)$. If the value of the function has moved with
$x$, we will therefore retrieve the value $I(x)$ at location $y =g(x)$
on the deformed image, $I'$. We therefore have $I'(g(x)) = I(x)$ or $I' =
I\circ g^{-1}$.

The transformation $I \mapsto I\circ g^{-1}$ represent the action of
the diffeomorphism $g$ on functions defined on $\Om$. The matching
problem in this case can be formulated as: given $I$ and $I'$, find a
diffeomorphism $g$, as ``small'' as possible, such that $I\circ g{-1}$
is similar to $I'$.
 

\subsection{Variational approaches}

Starting with \cite{bb82}, a very large number of algorithms have been
based on variational formulations. In this framework, one looks for
the minimum of
\begin{equation}
\label{eq:var.U}
U(g) = S(g) + \lambda D(I\circ g^{-1}, I')
\end{equation}
where $\lambda$ is a est un parameter, $S(g)$ measures the size of the
diffeomorphism (like in section \ref{size}) and
$D$ is some metric that evaluates the difference between $I\circ g^{-1}$ and
$I'$, for example
$$
D(I\circ g^{-1}, I') = \int_\Om \left|I\circ g^{-1} - I'\right|^2 dx
$$


\subsubsection{1D matching}

\nocite{ma97, pj87, pd95, sc78, gmk88, ggcv95}

One of the first situations in which these techniques have been
applied is for speech processing. If two acoustic signals are recorded
during a certain time, they need to be synchronized before safely
compariung them. Although current methods use a probabilistic approach
for this problem (with hidded Markov models), earlier attemps have
modeled distorsions due to variations in speech speed by a
diffeomorphism acting on the time variable. a  signal $I$ becoming
$I\circ g^{-1}$. The dynamic time warping technique has been
introduced for this purpose \cite{sc78}. It can be summarized as
follows. Assume that the signals are discretized, denoting  $I = (I_1,
\ldots, I_N)$ and $I' = (I'_1, \ldots, I'_{N'})$. The original algorithm
minimizes a function of the form 
$$
\De(I,I') = \sum_{k=1}^N d(I_k, I'_{\psi(k)})
$$
where $\psi(k)$ is the intex of the point which is homologous to
$I_k$ in $I'$. This term is the direct analogue of the second term in
\eqref{eq:var.U}. The smoothness penalty in \eqref{eq:var.U} is
replaced in DTW by a hard constraint on the differences  $\psi(k+1)
- \psi(k)$. The minimization is performed by dynamic programming
(which accounts for the D in DTW). Dynamic programming algorithms will
be described in appendix \ref{prog.dyn}.

DTW in general does not incorporate a diffeomorphic constraint. An
interesting contribution that accounts for this can be found in
\cite{pst98}. There two continuous signals $I$ and $I'$ are matched by
minimizing 
$$
U(g) = \arccos\left[\int_0^1 \sqrt{{g}_x} dx\right] + \la
\int_0^1(I(g^{-1}(x)) - I'(x))^2 dx.
$$
The unusual smoothness term is in fact, as we shall see later, a
geodesic distance on a particular group of diffeomorphisms.

One dimensional matching has direct implications for shape comparison,
when the compared shapes are paremetrized (say by arc-length). Indeed,
they are, in this case, functions from some interval to $\mR^2$, and
therefore belong in the very general class of images we are
considering here. The only issue is that curves withe different
lengths are defined ion different intervals, but this can be addressed
by rescaling all curves to have length 1, with is justifed since it provides
a scale invariant matching. Two curves $I,I': [0,1] \to \mR^2$ can then be matched by
minimizing 
$$
\De(I, I') = S(g) + \la \int_0^1 |I'(t) - I\circ g(t)|^2 dt,
$$
although most of the applications in shape comparison prefer working
with invariant signatures of the curve (like the curvature) rather
than curves as $R^2$-valued functions.

There is an extra complication to this framework when working with
shapes, because they are generally represented as closed curves. This
creates an undeterminacy in the arc-length representation, which comes
from the choice of the origin of the parameterization (the point for
which $s=0$). In most cases, this choice is arbitrary, so that a valid
matching procedure must also ensure that the origins in both curves
are chosen consistently, inducing an extra variable on which the
quality of the correspondence must be optimized. From a mathematical
point of view, this extra variable can be taken into account by
assuming that parameterizations are defined on the unit circle rather
than the unit interval. Modeling $g$ as a diffeomorphism of the unit
circle incorporates the change of origin. Unfortunately, this
variation has very little use for practical algorithms, because the
very efficient dynamic programming implementation require knowing the
origins which is therefore estimated separately (or via an exhaustive search). 


\subsubsection*{Examples}
\begin{itemize}
\item A generic method for matching plane curves would apply
\eqref{eq:var.U} to a geometric representation of the curve. For
example, assuming that all curves are defined on $[0,1]$, $I(s)$ can
be the Euclidean curvature of the curve (which is rotation and
translation invariant) or the angle of the tangent with the horizontal
(which is scale and translation invariant), or any other signature,
being given that the fewer derivatives the signature involves, the
more robust the method is likely to be. A typical energy in this
context can be
$$
U(g) = \int_0^1 |g_x|^2 dx + \la \int_0^1 (I \circ g^{-1}(x) - I'(x))
dx.
$$

\item A variant of this formulation \cite{cas92} can be given as
follows. Let $I$ and $J$ be parameterized with arc-length, $I$ having
length 1 and  $J$ length $\al$. Matching is searched as a bijection
$g: [0,\al]\to [0,1]$, that minimizes
$$
U(g) = \int_0^\al (\ka_I(g^{-1}(s)) - \ka_J(s))^2 ds + \la \int_0^\al
\left|\der{}{s}(I(g^{-1}(s)) - J)\right|^2 ds
$$
where $\ka_I$ and $\ka_J$ are the curvatures of $I$ an $J$.

Let's
rewite this energy in a form which is closer to the variatioal problem
we have been discussing. Introduce the unit tangents
 $T_I$ and $T_J$, and expand the second integral which becomes,
letting $h=g^{-1}$ and using the fact that $I$ and $J$ are
parameterinzed with arc-length,
$$
\al + \int_0^\al {h}_s^2 ds - 2\int_0^\al {h}_s
\scp{T_J}{T_I\circ h}(s) ds
$$
If $(s'\mapsto J(s'))$ is reparameterized by $s'\mapsto s'/\al$, and if
$\tilde h(s) = h(\al s)$, the previous formula becomes
$$
\al + \al^{-1}\int_0^1 \tilde h_s^2 ds - 2 \int_0^1 {\tilde h}_s
\scp{T_J}{T_I\circ \tilde h}(s) ds
$$
which is 
$$
\al + \al^{-1}\int_0^1 \tilde h_s^2 ds - 2 \int_0^1 
\scp{T_J\circ \tilde h^{-1}}{T_I}(s) ds.
$$
This can finally be rewritten as 
$$
\al  - 2 + \al^{-1}\int_0^1 \tilde h_s^2 ds + \int_0^1
|T_J \circ \tilde h^{-1} - T_I|^2 ds.
$$
Going back to the initial energy, we have to minimize
\begin{eqnarray*}
\tilde U(\tilde h) &=&  \la \al^{-1}\int_0^1 \tilde h_s^2 ds +
\int_0^\al (\ka_I(\tilde h(s/\al)) -
\ka_J)^2 ds \\
&+& \la \int_0^1
|T_J \circ \tilde h^{-1} - T_I|^2 ds.
\end{eqnarray*}
The first term is an elastic deformation energy for $\tilde h$. The
last two compare the curvature and unit tangents at homologous points.
 \nocite{ggcv95}

\item Some curve-matching energy have been proposed on an axiomatic
basis. The problem with such approaches is that different choices of
axioms lead to different forms of energy, so that they cannot be
expected to provide a universal answer, but nonetheless define
interesting families of matching problems.

One of these approaches \cite{bcgj95} let to constraints for which a
typical energy is
$$
U(g) = \int_0^1 |g_s - 1| ds + \la \int_0^1 (g_s + 1) |f(\ka_I\circ g) -
f(\ka_J)|ds
$$
with $f(\ka) = c\ka - \text{sign}(\ka) e^{-\al \ka}$. Another choice is
$$
U(g) = \int_0^1 |g_s - 1| ds + \la \int_0^1 |g_s \ka_I\circ g  - \ka_J|ds
$$

In \cite{ty00}, another set of axioms has been used, leading to
energies of the kind
$$
U(g) = \int_0^1 \sqrt{\dg} F(I, I'\circ g^{-1}) ds.
$$
Here, $U$ must be maximized in $g$; $F$ must be a symmetric, positive
function and satisfy the condition
$$
F(u, v) \leq \sqrt{F(u, u) F(v, v)}.
$$

\item One particular instance of the energy $U$ in \cite{ty00} is \cite{you98}
$$
U(g) = \frac{ \sqrt{2}}{2} \int_0^1
\sqrt{\dg\left(1 + \scp{T_I\circ g}{T_J}\right)}ds
$$
We will give later a justification of this energy in terms of
distances between plane curves.



\item Energies similar to \eqref{eq:var.U} can also be devised for 3D
curves. A simple choice  \cite{bgkm00} is to use the intrinsic
representation of section \ref{sec:3D.curve}, chapter \ref{surfaces}
and minimize, $I$ and $J$ being two curves with curvature $\ka_I$ and
$\ka_J$ and torsion $\tau_I$ and $\tau_J$, the energy
$$
U(g) = \int_0^1 g_s ds + \int_0^1 \big((\ka_I\circ g^{-1} - \ka_J)^2 +
(\tau_I\circ g^{-1} - \tau_J)^2\big)ds.
$$
\end{itemize}

Most of the energies described above can be minimized, in a very
efficient way, by dynamic programming (once a choice of the origins
has been made for closed curves). More details can be found in
appendix \ref{prog.dyn}.


\subsubsection{Matching higher dimensional images}

The main problem, when passing to high dimensional images (high
starting at 2 here) is of numerical nature. Feasible problem in 1D can
be challenging in 2D and simply impossible in 3D. Solving such
problems require either large scale computational power, or some
efforts in devising suitable simplifications and approximations.

Let us first review some specific domains of application of image matching.

\subsubsection*{Optical flow and stereo-vision}
These two old problems in image processing involve in one way or
another the use of a matching procedure.

For motion estimation, a displacement field  $u$ (the optical flow)
must be estimated between two consecutive images  $I$ and $I'$, in
order that $I(x-u) \simeq
I'(x)$. There is a very rich literature on this subject, which is one
of the most important problems in image processing. We will not try to
provide an exhaustive account of the methods, which fall outside the
scope of this book anyway \cite{fp02}.

A large variety of variational approaches used for optical flow,
however, are variants of the one initially introduced by Horn and
Schunk  \cite{hs81}, in which one minimizes the energy  (with $g = \id - u$),
\begin{equation}
\label{hs.en}
U(g) = \int_{\Om} \al^2 |u(x)|^2du +  \int_{\Om} |I(g(x)) -
I'(x)|^2 dx.
\end{equation}
Most of the time, the second term is lenearized in a neighborhood of
$u=0$ and replaced by
$$
I(x) - I'(x) - \nabla I(x).u(x)
$$

In stereo vision \cite{msks03}, matching two images that have been
simultaneously acquired by two cameras allows one to compute a depth
map within the matched parts of the two images. Compared to the
optical flow, there is an additional constraint (called epipolar)
which ensures that homologous points belong to specific lines which
only depend on the relative position of the cameras and their internal
geometry. 


\subsubsection*{Face comparison}
\nocite{bhk97} 

Deformation also plays an important role in the analysis of facial
images. One can list three main factors for the variation of face images
of a single person: illumination, position and expression. We focus on
the latter than can be represented as a diffeomorphism (for a complete
theory including the three factors, we refer to 
\cite{hgygm99}).

Although using non-rigid matching methods is not systematic in this
domain, some fundamental contributions have been made using this
technique \cite{fe73}. In the work of  Von
der Malsburg and his collaborators, for example \cite{dkl92, wv93},
face comparison incorporates a discrete form of elastic deformation,
represented as neural networks; the aligned quantities are not image
grey levels, but Gabor filters \cite{gab46}, in accordance with biological
observations.

In \cite{hgygm99}, a (quite complex) model is proposed for face variation in terms of illumination,
position and expression, to provide a single variational problem. The
deformation part incorporates the hyperelastic energy we have
described in equation
(\ref{elas.ax}).




\subsubsection*{Medical imaging}

Medical imaging may be the primary domain of application of
deformations, in 2, but mainly 3, dimensions.
Increasingly sophisticated acquisition systems provide anatomical
images in higher number with more and more accuracy, with new informations and
modalities. The need to, at least partially, automatize the processing
and analysis of these images has become more and more important. The
whole field of computerized analysis of medical images is called
computational anatomy \cite{gm98}. Several matching methods applied to
brain images are described in  \cite{tog99}.

\subsection{Using flows of diffeomophisms}
\label{fluid.grad}

Like for interpolation, the previous techniques do not 
guarantee that the optimal deformation is a diffeomorphim. Such a
property can be obtained by using the formulation given in chapter
\ref{grp.diff}. This leads to minimizing
$$
U(\bfv) = \int_0^1 \int_\Om \scp{L\bfv(t, .)}{\bfv(t, .)} dx dt + \la \int_\Om
\left(I\circ \bfg_\bfv(1,x) - I'(x)\right)^2 dx
$$
with respect to  $\bfv$, where $L$ is a strictly monotonous admissible
operator (i.e., such that $\CH_L$, according to definition \ref{adm}),
and, as before,  $\bfg_\bfv(t,x)$ is the value at time 
$t$ of the solution of $y_t = \bfv(t, y)$ with initial condition
 $y(0) = x$ ($x\in\Om$).

Because the operator $L$ is admissible, the function  $U$ achieves a
global minimum, for a time-dependent vector filed  $\bfv$ that
generates a diffeomorphim. Let's prove this in the case when  $I$ and
$I'$ are continuous on $\Om$ (the result remains true for $L^2$
functions). The mapping
$$
g\mapsto q(g) = \int_\Om |I'(x) - I\circ g(x)|^2 dx
$$
is continuous for uniform convergence: if $g_n$ converges to $g$
uniformly on $\Om$, then, from the continuity of $I$ and the dominated
convergence theorem,  $q(g_n)$ converges to $q(g)$. This fact,
combined with theorem \ref{weak.closed} and to the lower
semi-continuity of norms on Hilbert spaces implies that there exists a
$\bfv^*$ minimizing $U(\bfv)$.

\bigskip
To minimize this expression using gradient descent, we need to compute
the differential of of $U$ with respect to  $\bfv$. This require some
care, in particular for the term  $\bfg_\bfv$. We now proceed to this
computation assuming that $I$ and $I'$ are $C^1$ functions. We
evaluate the 
variation $U(\bfv+ \ep h) - U(\bfv)$ for small $\ep$.

The first term of $U(\bfv+\ep h) - U(\bfv)$ is $\|\bfv+\ep h\|^2 -
\|\bfv\|^2$, which can be written $\ep \scp{h}{\bfv} + o(\ep)$ with
$$
\scp{h}{\bfv} = \int_0^1 \int_\Om \scp{Lh}{\bfv} dx dt.
$$
The study of the second term, which is
$$
\int_\Om \left(I\circ \bfg_{\bfv+\ep h}(1,x) - I'(x)\right)^2 dx -
\int_\Om \left(I\circ \bfg_{\bfv}(1,x) - I'(x)\right)^2 dx
$$
requires more work.

We first show that the differential
$$
q_h(t) = \pdr{}{\ep}\bfg_{\bfv+\ep h}(t, x)_{|_{\ep=0}}
$$
exists. This can be handled in the same way as the existence of the
differential  of $\bfg_\bfv$ with respect to $x$ (theorem
\ref{flot.exist}), and we rapidly sketch the details. One first
formally computes a candidate by writing
$$
\pdr{}{t} \bfg_{\bfv+\ep h}(t, x) = \bfv(t, \bfg_{\bfv+\ep h}(t, x)) +
\ep h(t, \bfg_{\bfv+\ep h}(t, x))
$$
to get
\begin{equation}
\label{qh}
\der{q_h}{t} = D \bfv(t, .) (\bfg_\bfv(t, x)) q_h + h(t,\bfg_{\bfv}(t,x)).
\end{equation}
This is a linear, non-homogeneous differential equation, for which we
introduce a solution $\tilde{q}_h$ (with
$\tilde{q}_h(0) = 0$). It remains to make this formal computation
rigorous, by showing that 
$q_h$ does exist, and is equal to $\tilde{q}_h$. This can be done by
returning to the definition of the derivative and use Gronwall's lemma
like in the proof of theorem \ref{flot.exist}.

Remarking that the homogeneous equation associated to \eqref{qh} is in
fact \eqref{flot.diff},  $q_h$ can be computed analytically by the
mathod of variation of the constant, yielding
$$
q_h(t, x) = D\bfg_\bfv(t,x)\int_0^t D\bfg_\bfv(s,x)^{-1}.h(s,
\bfg_\bfv(s,x)) ds.
$$

The previous computation therefore yields
\begin{multline*}
U(\bfv+\ep h) - U(\bfv) = 2 \ep \scp{h}{\bfv} \\
+ 2\ep \la \int_\Om \scp{I\circ
\bfg_\bfv(1,x) - I'(x)}{DI\circ \bfg_\bfv(1,x)q_h(1,x)} dx
+
o(\|h\|).
\end{multline*}
Replacing  $q_h$ b its expression, we get, denoting 
$\ga_\bfv(t, .)$ the inverse of $\bfg_\bfv(t,.)$
\begin{multline*}
\int_\Om \scp{I\circ \bfg_\bfv(1,x) - I'(x)}{DI\circ
\bfg_\bfv(1,x).q_h(1,x)} dx
\\ = \int_0^1\int_\Om \left\langle I\circ
\bfg_\bfv(1,x) - I'(x), \right. \\
\left.DI\circ
\bfg_\bfv(1,x).d\bfg_\bfv(1,x)d\bfg_\bfv(t,x)^{-1}.h(t,
\bfg_\bfv(t,x)) \right \rangle dx dt
\\=
 \int_0^1\int_\Om \lan  DI\circ
\bfg_\bfv(1,\ga_\bfv(t,y))D\bfg_\bfv(1,\ga_\bfv(t,y))D\bfg_\bfv(t,\ga_\bfv(t,y))^{-1}.
h(t, y),
\\ I\circ \bfg_\bfv(1,\ga_\bfv(t,y)) -
I'(\ga_\bfv(t,y))\ran |\det D\ga_\bfv(t,y)|dy dt.
\end{multline*}
To simplify this formula, introduce the function $\chi_\bfv(t, y) = \bfg_\bfv(1,
\ga_\bfv(t,y))$, and the notation $\bfJ(t,y) = I(\chi_\bfv(t,y))$
and $\bfJ'(t,y) = I'(\ga_\bfv(t,y))$. We have
$$D\bfg_\bfv(t,\ga_\bfv(t,y))^{-1} = D\ga_\bfv(t,y)\,,
\ \ 
D\bfg_\bfv(1,\ga_\bfv(t,y)) D\ga_\bfv(t,y) = D\chi_\bfv(t,y)
$$
and
$\dps
D\bfJ(t,y) = DI\circ \bfg_\bfv(1,\ga_\bfv(t,y)) D\chi_\bfv(t,y)\,.$
We finally obtain the expression
\begin{multline*}
\int_\Om \scp{I\circ \bfg_\bfv(1,x) - I'(x)}{DI\circ
\bfg_\bfv(1,x) q_h(1,x)}
\\ =   \int_0^1\int_\Om \scp{\bfJ(t,y) - \bfJ'(t,y)}{D\bfJ(t,y) |\det
D\ga_\bfv(t,y)| h(t, y)} dy dt \\ = \int_0^1\int_\Om |\det
D\ga_\bfv(t,y)|\scp{D\bfJ(t,y)^T (\bfJ(t,y) - \bfJ'(t,y))}{h(t,
y)} dy dt
\end{multline*}
and the proposition
\encadre{
\begin{proposition} [Beg et al., 2002]
\label{grad.prop}
Let
$$
U(\bfv) = \int_0^1 \int_\Om |L\bfv|^2 dx dt + \la \int_\Om
\left(I\circ \bfg_\bfv(1,x) - I'(x)\right)^2 dx
$$
If $u \mapsto \left(\int_\Om |Lu|^2 dx\right)^{1/2}$ induces an
admissible norm on functions $u:\Om\to\mR$, we have, letting
$\|\bfv\|^2 = \int_0^1 \int_\Om |L\bfv|^2 dx dt\,:$
$$
U(\bfv+\ep h) = 2 \ep \scp{\bfv}{h} + 2\ep \la \int_0^1\int_\Om
\scp{b_\bfv(t,y)}{h(t,y)} dy + o(\ep)
$$
with
$$
b_\bfv(t, y) = |\det D\ga_\bfv(t,y)|D\bfJ(t,y)^T (\bfJ(t,y) -
\bfJ'(t,y))
$$
where $\ga_\bfv(t, y)$ is the inverse of $\bfg_\bfv(t,y)$ and
$
\bfJ(t, y) = I(\bfg_\bfv(1,\ga_\bfv(t, y)))$, $\bfJ'(t,y) =
I'(\ga_\bfv(t, y))$.
\end{proposition}
}

Figure \ref{match.beg} provides a matching experiment using this
algorithm for shape matching.

\begin{figure}
\begin{center}
\includegraphics[width=5.5cm]{TUTORIALPS/SY0.ps}\hfill
\includegraphics[width=5.5cm]{TUTORIALPS/SY299.ps}\\
\includegraphics[width=5.5cm]{TUTORIALPS/SYgrid270.ps}
\end{center}
\caption{\label{match.beg} First Row: two shapes before
alignment. Second Row: Estimated deformation applied to a regular grid.}
\end{figure}


\subsection{Numerical aspects}
In this section, we make (without geting into imlpementation details)
some important remarks on the numerical 
implementation of diffeomorphic matching. 

\subsubsection{Discretization}
The simplest and mostly used way to discretize a function is to use a
regular grid, approximate integrals by sums and derivatives by finite
differences. This leads to rather simple implementations, but involves
a huge number of variables, that can require considerable computer
power in 3D.

One can also use a decomposition on a basis of functions. One
possibility is to use finite elements 
\cite{rwcm95}, which consists  \cite{zt89} in defining a basis of
smooth function (generally piecewise polynomial, with compact support)
on which the unknow diffeomorphism is projected. Another possible
choice of basis are the eigenvectors of the operator $L$ which form an
orthonormal basis of  $\CL^2$.


\subsubsection{Definition of the gradient}

Here is a simple, but fundamental remark, on gradient descent
algorithms. At the difference of the differential of a function, the
gradient depens on a dot product (or a Riemannian structure) on the
ambient space. For example, if $U$ is a differentiable function on
$\mR^p$, the differential of $U$ at $g$ is a linear operator, $DU(g)$,
defined by 
$U(g+\ep h) \simeq U(g) + \ep DU(g).h + o(\ep)$, avec
$$
DU(g).h = \sum_{i=1}^p \pdr{U}{g_i} h_i.
$$
To define a gradient, we need a Eclidean or Riemannian
structure. Denoting it  $\scp{.}{.}_g$, the gradient of  $U$ at $g$ is {\em defined}
by
$$
\forall h,  DU(g).h = \scp{\nabla U(g)}{h}_g.
$$
The discrete gradient descent algorithm with step  $\ga$ is
$$
g(\tau+1) - g(\tau) = -\ga \nabla U(g(\tau)).
$$
This algorithm therefore also depends on the Riemannian structure. On
$\mR^p$, the standard Euclidean structure provides the usual gradient
descent, 
$$
g(\tau+1) - g(\tau) = - \ga \der{U}{g}
$$
where $\der{U}{g}$ is the vector composed with partial derivatives of $U$
at $g(\tau)$; but, letting
$$
\scp{u}{v}_g = {}^tuA_gv,
$$
where $A_g$ is a symmetric positive matrix, the algorithm becomes
$$
g(\tau+1) - g(\tau) = - \ga A_{g(\tau)}^{-1}\der{U}{g}.
$$

Such variations may have importance consequences on the quality and
stability of the algorithm, especially in infinite dimensions. From a
heuristic point of view, the Riemannian metric defines directions that
are easier to take than other in the space, and make others almost
impossible. If  $g$ is a function, a suitable choice of the metric
will favor trajectories where $g$ remains bounded or smooth, for example.
Even in finite dimensions, selecting a good metric can lead to
improved algorithms. For example, when a function is optimized over
invertible matrices, one can ensure that the matrix remains invertible
along the algorithm (like for independent component analysis, appendix
\ref{ica}). 

Let's provide some illustrations of this for the optimization over
functions. Consider the following typical situation in the calculus of
variation where a function $g$ defined n $\O$ is searched to minimize
$$
U(g) = \int_\Om F(Dg(x), g(x), x) dx.
$$
We first compute the equivalent of directional derivatives, which
coreespond to computing a function $b_g$ defined on $\Om$, such that,
for every function $h$ with compact support in $\Om$
$$
\der{}{\ep} U(g+\ep h) = \int_\Om b_g(x).h(x) dx
$$
at $\ep=0$. This can be interpreted as a gradient of $U$ for the
$\CL^2$ dot product
$$
\scp{h}{k}_2 = \int_\Om k h dx\,.
$$

Take, for example, $\Om = ]0,1[$ and
$$
U(g) = \int_0^1 (|\dg|^2 + V(g(x))) dx
$$
where $V$ is a positive function and  $g: [0,1]\to [0,1]$ satisfies
$g(0) = 0$ and $g(1) = 1$. By a standard computation using integration
by parts, we obtain, at $\ep=0$,
$$
\der{}{\ep}U(g+\ep h) = \int_0^1(-\dder{g}{x} + V'\circ g).h dx.
$$
This provides the $\CL^2$-gradient of $U$
$$
\nabla_g U = -\dder{g}{x} + V'\circ g.
$$
The gradient dscent algorithm (expressed in continuous time) is
$$
\pdr{g}{\tau} = -\dder{g}{x} + V'\circ g.
$$

\bigskip

Return now to the general case. Assume that the dot product is
selected so that 
$$
\scp{h}{k} = \int_\Om \scp{k}{Lh} dx
$$
where $L$ is a strictly monotnus operator. Still assuming that we can
write
$$
\der{}{\ep}U(g+\ep h) = \int_\Om \scp{b_g}{h} dx\,,
$$
we see that the gradient of $U$ for the new dot.product must satisfy
$Lv = b_g$, yielding the numerical scheme
$$
\left\{
\begin{array}{l}
\dps \pdr{g}{\tau} = - v\\
\ \\
\dps Lv = b_g
\end{array}
\right. .
$$
In this discussion, we have neglected the boundary conditions on $\prt
\Om$ (or implicitely assumed that valid variations $h$ must vanish at
the boundary, like in our example, so that they do not add any new
term in the integration by parts).

Let's return to out 1D example, and consider the dot product
$$
\scp{k}{h}_L = \int_0^1 \dot{k}\dh dx = \int_0^1 k Lh dx
$$
with $Lh = - \dder{h}{x}$; this yields the algorithm
$$
\left\{
\begin{array}{l}
\dps \pdr{g}{\tau} = - v\\
\ \\
\dps \dder{v}{x} = -\dder{g}{x} + V'\circ g\\
\ \\
\dps v(0) = v(1) = 0
\end{array}
\right.
$$

We have therfore built a two-step numerical scheme. First compute the
$\CL^2$ gradient then {\em smooth it}, by applying to it the inverse
of the operator $L$, to obtain the increment of the crent function
$g$. More precisely, if $K = L^{-1}$ is a kernel operator, the
algorithm is
$$
\pdr{g}{\tau} = -\int_{\Om} K(x, y) b_g(y) dy.
$$
This implies in particular that $g$ inherits the smoothness of the
kernel during evolution, which is generally a factor of stability in
the algorithm.

The situation is even more interesting in the case of groups of
diffemorphisms, in which we deal here with a ``true'' Riemannian metric and
not a change of Hilbert space dot product. The application of this
idea to groups of diffeomorphisms has been suggested in
\cite{tro98}. Consider the function
\begin{equation}
\label{eq:image.data}
U(g) = \int_\Om \left| I \circ g(x) - J(x) \right|^2 dx
\end{equation}
where $I$ and $J$ are two $C^1$ images defined on $\Om$, and $g$ is a
diffeomorphism of $\Om$. It is very easy to compute the
$\CL^2$ gradient of $U$, and is given by
\begin{equation}
\label{grad.l2}
b_g(x) = 2(I(g(x)) - J(x)) DI(g(x)).
\end{equation}
The $\CL^2$ gradient descent would therefore be given by
$$
\pdr{g}{t} = - (I(g(x)) - J(x)) DI(g(x)).
$$

This procedure is extremely unstable. It is in fact never used as such,
and always combined with some form of regularization (see the method
of deamons in the next section). Another issue is that the procedure
does not take into acount the fact that  $g$ must be a
difeomorphism. Let's consider a Riemannian dot product for
diffeomorphisms, given by
$$
\scp{h}{k}_g = \int_\Om \scp{L(h\circ g^{-1})}{k\circ g^{-1}} dx.
$$
(The operator $L$ is applied to the function $h\circ g^{-1}$.) This
choice will be justfified in chapter \ref{dist.diff}. To obtain the
associated gradient descent, we need to express
\begin{equation}
\label{eq:grad.img}
\der{}{\ep}U(g+\ep h)  \int_\Om 2(I(g(x)) - J(x)) DI({g(x)}) h(x)
dx 
\end{equation}
under the form
$$
\der{}{\ep}U(g+\ep h) = \scp{\beta_g}{h}_g. $$
Let us first make the change of variables  $y = g(x)$ in
\eqref{eq:grad.img}. This yields
$$
\der{}{\ep}U(g+\ep h) = \int_\Om 2(I(y) - J(g^{-1}(y)))
DI(y) h(g^{-1}(y)) |\det Dg^{-1}(y)| dy + o(\ep).
$$
This implies that  $\beta_g$ must be such that $\beta_g\circ g^{-1}$
satisfies the equation
$$
L.u(y) = 2|\det Dg^{-1}(y)|(I(y) - J(g^{-1}(y)))
DI(y).
$$

This means that the gradient descent for the considered dot product
must proceed as follows.Solve
\begin{equation}
\label{eq:flow.grad}
\pdr{g}{\tau}(\tau,x) = u(\tau, g(\tau,x))
\end{equation}
with 
\begin{equation}
\label{eq:smooth.grad}
L(t,.) = - 2|\det D{g}^{-1}(y)|(I(y) - J({g}^{-1}(y)))
DI(y).
\end{equation}
Equation \eqref{eq:smooth.grad} is a smoothing operation applied to
the $\CL^2$ gradient. The particular form of the evolution equation
\eqref{eq:flow.grad} is a dircet consequence of the form that has been
selected for the Riemannian metric. It shows that the evolution
generated by this particular
implementation of gradient descent s a flow (a solution of an ODE)
that generates a diffeomorphism as long as $u$ remains smooth. A
similar algorithm as been proposed by \cite{crm96}, based on different
motivations. 

It is important to note that that the function $U$ in \eqref{eq:image.data} contains no
regularization term for $g$. This implies that, unless $I$ and $J$ are
can be made in correspondence exactly by diffeomorphism, there is
probably no global minimum, and the evolution will generate wild
deformations if pushed too far. A suitable stopping rule must be
designed \cite{tro98}. 




\subsubsection{Deamons algorithm}

The method of deamons is a popular matching algorithm in medical
imaging \cite{thi98, thi99, gram99}. Although it does not exactly
correspond to a gradient descent implementation, it can almost be
interpreted as such. Consider the energy
$$
U(g) = \int_\Om |J(x) - I(g(x))|^2 dx
$$
with $\CL^2$ gradient given by
$$b_g(x) = 2(I(g(x)) - J(x)) DI ({g(x)})$$

Consider the following dot product
$$
\scp{h}{k}_g = \int_\Om (|J(x) - I(g(x))|^2 + |D J(x)|^2) h(x)
Lk(x) dx
$$
for some operator $L$.
The gradient of  $U$ for this dot product is given by
$$
L^{-1} \Big(\frac{b_g(x)}{|J(x) - I(g(x))|^2 + |D J(x)|^2}\Big)
$$
that can be approached by
$$
-2L^{-1}\left(\frac{J(x) - I(x)}{|J(x) - I(g(x))|^2 + |D J(x)|^2} D
J(x)\right).
$$

This provides a gradient descent algorithm which can be written, in
discrete time, under the form (letting $K=L^{-1}$)
\begin{equation}
\label{thi}
u(\tau+1, y) = u(\tau, y) + \ga\int K(y,x) \frac{J(x) - I(x)}{|J(x) - I(g(\tau, x))|^2 + |\nabla J(x)|^2} \nabla
J(x) dx
\end{equation}
with $g(\tau, x) = x + u(\tau, x)$. The method of deamons applies an
additional smoothing to the current $u$ before adding the increment,
yielding
\begin{multline*}
u(\tau+1, y) = \int K(y,x) u(\tau,x))dx\\
 + \int K(y,x) \ga\frac{J(x) - I(x)}{|J(x) - I(g(\tau, x))|^2 + |\nabla J(x)|^2} \nabla
J(x) dx.
\end{multline*}
This addictional smoothing is hard to interpret and to justify.
To interpret the dot product introduced in this section, one must
remember that the norm
$\scp{h}{h}_g$ can be interpreted as an infinitesimal cost function
for the variation $g\to g+h$; as we have already noticed, the operator
$L$ penalizes a lack of smoothness; the multiplicative factor
penalizes variations at location $x$ according to the value of  $|J(x)
- I(g(x))|^2 + |D J(x)|^2$. This implies that changes of $h(x)$
are limited when the matching is unsatisfactory and when $J$ has a
large gradient. 



\nocite{yos70}

\subsubsection{Multigrid approaches}

For these very high dimensional, non-convex, multigrid algorithms
often bring significant improvements, both in efficiency and
quality, in avoiding some local minima. Several strategies are available for this.
One can start solving a problem with coarser images and discretization grids, to
progressively refine the solutions. One can also use representations
of the unknown variable (the diffeomorphism) that allow for a
progressive release of the degrees of freedom. This can be done easily
by using a decomposition on orthogonal bases and playing with the
number of coefficients that are retained \cite{bk89, ghlb97, bk99}.









\chapter{Distances and group actions}
\label{dist.grp}

\section{General principles}
\subsection{Distance induced by a group action}

This chapter shows how transformation groups acting on spaces can help
defining or altering distances on these spaces. The first part will
essentially focus of the algebraic aspects of the contructions based
on the least action principle. In the second part, we will develop the
related
differential point of view, when a Lie group acts on a manifold.

\index{distance (axioms)}

\bigskip
We start with some terminology. A distance on a set  $\CC$ is a mapping
$d:\CC^2\mapsto \mR_+$ such that, for all $m, m', m''\in\CC$,\\
\indent D1) $d(m, m') = 0 \Leftrightarrow m= m'$.\\
\indent D2) $d(m. m') = d(m', m)$\\
\indent D3) $d(m, m'') \leq d(m, m') + d(m', m'')$.

If D1 is not satisfied, but only the fact that $d(m, m)=0$ pour tout
$m$, we will say (still assuming D2 and D3) that $d$ is a {\em pseudo-distance}.

If $G$ is a group acting on $\CC$, we will say that a distance $d$ on
$\CC$ is $G$-equivariant if and only if for all $g\in G$, for all
$m, m'\in \CC$, $d(g.m, g.m') = d(m, m')$.

A mapping $d:\CC^2\mapsto \mR_+$ is a $G$-invariant distance if and
only if it is a pseudo-distance such that $d(m, m') = 0
\Leftrightarrow \exists g\in G,\, g.m = m'$. This is equivalent to
stating that $d$ is a distance distance on the coset space $\CC/G$,
composed with orbits
$$
[m] = \{g.m, g\in G\},
$$
with the identification $d([m], [m']) = d(m,, m')$.

The next proposition shows how a $G$-equivariant distance induces a
$G$-invariant pseudo-distance.
%\encadre{
\begin{proposition}
\label{modulo}
Let $d$ be equivariant for the left action of $G$ on $\CC$. The function $\td$,
defined by
$$
\td([m], [m']) = \inf\{d(gm, g'm'), g, g'\in G\}
$$
is a pseudo-distance on $\CC/G$.
\end{proposition}
%}
\index{invariant distance} \index{equivariant distance}
Since $d$ is $G$-equivariant, we could have defined $\td$ in the
previous proposition by
$$
\td([m], [m']) = \inf\{d(gm, m'), g\in G\}.
$$
\begin{proof}
The symmetry of $\td$ is obvious, as well as the fact that $\td((m],
[m]) =0$ for all $m$. For the triangle inequality, D3, it suffices to
show that, for all $g_1, g'_1, g'_2, g''_1\in G$, there exists $g_2,
g''_2\in G$ such that
\begin{equation}
\label{eq:d3}
d(g_2m, g''_2 m'') \leq d(g_1m, g'_1m') + d(g'_2m', g''_1m'').
\end{equation}
Indeed, if this is true, the minimum of the right-hand term in $g_1, g'_1, g'_2, g''_1$, which is
$\td([m], [m']) + \td([m'], [m''])$ is larger than the minimum of the
left-hand term in $g_2, g''_2$, which is $\td([m], [m''])$, which is
precisely D3. To prove \eqref{eq:d3}, write $d(g'_2m', g''_1m'') = d(g'_1m', g'_1 (g'_2)^{-1}
g''_1m'')$, and take $g_2 = g_1$ and 
$g''_2 = g'_1 (g'_2)^{-1}
g''_1$; \eqref{eq:d3} is just an application of the triangle
inequality for $d$.
\end{proof}

Whether $\td$ is a distance, and not just a pseudo-distance is
statement about the minimizer of a function. Excepted in the trivial
case of finite groups, this will require some topological assumptions
on the structure of $G$ and of the action.

The same statement can clearly be made with $G$ acting on the right on $m$,
with $m \mapsto m.g$. We state it without proof
\begin{proposition}
\label{modulo.right}
Let $d$ be equivariant for the left action of $G$ on $\CC$. The function $\td$,
defined by
$$
\td([m], [m']) = \inf\{d(mg, m'g'), g, g'\in G\}
$$
is a pseudo-distance on $G\backslash \CC$.
\end{proposition}
Here $G\backslash M$ denotes the coset space for the right action of $G$.



\subsection{Distance altered by a group action}

In this section, $G$ is still a group acting on the left on  $\CC$, but we consider
the product space $\CO  = G\ti \CC$ and project on $\CC$ a distance
defined on $\CO$.

The group $G$ acts on the right on $\CO$, letting, for $k\in G$,
$o=(h, m)\in \CO$, $o.k = (hk, k^{-1}m)$. For $o=(h, m)\in \CO$, we define
the projection 
$\pi(o) = hm$. This projection is such that, for all $k\in G$,
$\pi(o.k) = \pi(o)$. 

Let $d_\CO$ be a distance on $\CO$. We let, for $m, m'\in \CC$
\begin{equation}
\label{inv.eq.gen}
d(m, m') = \inf\{d_\CO(o, o'), o, o'\in \CO, \pi(o) = m,  \pi(o') = m'\}.
\end{equation}

We have the proposition
\encadre{
\begin{proposition}
\label{inv.gch.gen}
If $d_\CO$ is equivariante by the right-action of $G$, then, the
function $d$ defined by \ref{inv.eq.gen} is a pseudo-distance on $\CC$.
\end{proposition}}

This is in fact a corollary of proposition \ref{modulo.right}. One only has
to remark that the quotient space  $G\backslash \CO$ can be identified to $\CC$
via the projection $\pi$, and that the distance in \eqref{inv.eq.gen}
then becomes the projection distance introduced in proposition
\ref{modulo.right}.

\subsection{Existence of a minimum}

We assume that the group  $G$ is a topological space, but no special
structure on $\CC$.
For $m\in \CC$, define the mapping $I_m: G\mapsto \CO$ by
$I_m(h) = (h, h^{-1}m)$. We make the following hypotheses.
\begin{itemize}
\item The canonical projection  $(h, m)\in
\CO \to h$ of sets of finite diameters for  $d_\CO$ are sequentially
precompact sets in $G$.
\item For all $o\in \CO$, for all $m\in \CC$, the function $h\mapsto
d_\CO(o, I_m(h))$ is lower semi-continuous: if $h_n\to
h$ in $G$, then
$$
d_\CO(o, I_m(h)) \leq \liminf d_\CO(o, I_m(h_n))
$$
\end{itemize}
A subset  $A$ of the topological space $G$ is sequentially precompact
if one can extract a sequence that converges in $G$ from any sequence
in  $A$.

We have
\begin{proposition}
\label{exist.gen}
Assume the previous hypotheses. If $d_\CO$ is equivariant by the
left-action of  $G$, then, for all  $m,
m'\in\CC$, there exists $o, o'\in G$ such that $\pi(o) = m$,
$\pi(o')=m'$ and $d(m, m') = d_\CO(o, o')$. In particular, $d$ is a
distance on  $\CC$.
\end{proposition}
\begin{proof}
Let $m$ and $m'$ be two elements of $\CC$. From the definition of
$d$, there exists two sequences  $(o_n, n\geq 0)$ and $(o'_n, n\geq
0)$ such that  $\pi(o_n) = m$, $\pi(o'_n) = m'$, and $d_\CO(o_n, o'_n) \ria
d(m, m')$. From the definition of  $\pi$, $o_n$ takes the form $o_n =
(h_n, h_n^{-1}.m)$, and $o'_n = (h'_n, {h'_n}^{-1}.m')$. Using the
left-equivariance, we can replace  $o_n$ and $o'_n$ by $h_n o_n$ and
$h_n o'_n$; this is equivalent to assuming that  $h_n = \id_G$ for
all $n$, that we shall do without changing the notation, letting
$o_n \equiv o=(\id_G, m)$ for al $n$. Since $o_n$ is constant, the
sequence  $(o'_n, n\geq 0)$ is bounded in $\CO$, so that we can
extract a converging subsequence from  $(h'_n, n\geq 0)$; we let
$h'\in G$ be the limit. We have $o'_n = I_{m'}(h'_n)$, and, letting
$o'=(h', h'^{-1}m') = I_{m'}(h')$, our second hypothesis yields
$$
d_\CO(o, o') \leq \liminf d_\CO(o, o'_n) = d(m, m').
$$
Since $d(m, m') \leq d_\CO(o, o')$, we obtain the required result.
\end{proof}

\subsection{Transitive action}
\subsubsection{Induced distance}
An important particular instance of the previous construction is when
$\CO$ itself is a group that acts transitively on $\CC$. This means
that for any $m, m'$ in $\CC$, there exists an element $o\in\CO$ such
that $m' = om$.

So our initial assumption is that $\CO$ is a group that acts
transitively on the set $\CC$. We fix a reference element 
 $m_0$ in $\CC$, and define the group $G$ by
$$G = \{o\in \CO, o.m_0 = m_0\}\,.$$
This group $G$ is called the stabilizer of  $m_0$ in $\CO$. The
relation to the previous section comes with the identification of
$\CO$ with $G\ti\CC$. 

One way to see that is by assuming that a function  $\rho: \CC\to \CO$
has been defined, such that for all $m\in\CC$, $m = \rho(m).m_0$. This
is possible, since the action is transitive (using anyway
the axiom of choice). Define
$$
\map{\Psi}{G\ti\CC}{\CO}{(h, m)}{\rho(h.m)h}.
$$
We show that $\Psi$ is a bijection: if $o\in \CO$, one can write, we
look for $(h,m)$ such that 
$o = \Psi(h, m)$. If this is true, we have 
$$o.m_0 = \rho(hm)h m_0 = \rho(hm)m_0 = hm$$
which imples that $\rho(hm) = \rho(om_0)$ and therefore $h =
\rho(om_0)^{-1} o$ which is uniquely specified; but this also specifies $m =
h^{-1}om_0$. This proves that $\Psi$ is one-to-one and onto
and provides the identification we are looking for.

The right action of $G$ on $G\ti\CC$, which is $(h, m)k = (hk,
k^{-1}m)$, translates to $\CO$ via $\Psi$ with
$$
\Psi(hk, k^{-1}m) = \rho(hkk^{-1}m)hk = \Psi(h,m)k
$$
which coincides with the right action of $G$ on $\CO$. Finally, the
constraint $\pi(h,m_1) = m$ which is involved in 
Proposition \ref{inv.gch.gen} becomes $om_0 = m$ via the
identification. All this provides a new version of Proposition
\ref{inv.gch.gen} in our present context given by
\begin{corollary}
\label{inv.gch}
Let $d_\CO$ be a distance on $\CO$, which is invariant under the
right-action of  $G$, the stabilizer of $m_0\in\CC$. Define, for all
$m, m'\in\CC$
\begin{equation}
\label{inv.eqg}
d(m, m') = \inf\{d_\CO(o, o'), o.m_0 = m, o'.m_0=m'\}\,.
\end{equation}
Then $d$ is a pseudo distance on  $\CC$.
\end{corollary}

The translation of the hypotheses made for 
proposition  \ref{exist.gen} is as follows. We assume that  $G$ can
be equipped with a topology such that
\begin{itemize}
\item Images of sets of finite diameters for  $d_\CO$ by $o \mapsto
\rho(om_0)^{-1} o$ are sequentially compact sets in  $G$.
\item For all  $o\in \CO$, for all $m\in \CC$, the function  $h\mapsto
d_\CO(o, \rho(m)h)$ is lower semi-continuous..
\end{itemize}
We have
\begin{corollary}
\label{exist.gch}
We assume the previous hypotheses. If $d_\CO$ is equivariant by the
right action of  $G$, then, for all $m, m'\in\CC$, there exists $o,
o'\in \CO$ such that $\pi(o) = m$, $\pi(o')=m'$ and $d(m, m') = d_\CO(o,
o')$. In particular, $d$ is a distance on $\CC$.
\end{corollary}





\subsubsection{Effort functional}
\bigskip
\index{Effort functional}
As formalized in \cite{gre93}, one can built a distance on
 $\CC$ on which a group acts transitively using the notion of effort
functionals. The definition we give here is slightly more general than
in  \cite{gre93}, to take into account a possible influence of the
deformed object on the effort. We also make a connection with the
previous, distance based, formulations.

We attribute a cost $\Ga(o, m)$ to  a transformation $m\ria o.m$. If
$m$ and 
$m'$ are two objects, we define 
$d(m, m')$ as the minimal cost (effort) required to transform  $m$ in
$m'$, i.e.,
\begin{equation}
\label{ga.gch}
d(m, m') = \inf\{\Ga(o,m), g\in G, o^{-1}.m = m'\}
\end{equation}

The proof of the following proposition is almost obvious.
\begin{proposition}
\label{cout.gch}
If $\Ga$ satisfies

C1) $\Ga(o, m) = 0 \Leftrightarrow o=e$

C2) $\Ga(o, m) = \Ga(o^{-1}, om)$

C3) $\Ga(oo', m) \leq \Ga(o, m) + \Ga(o', om)$

Then, $d$ defined by (\ref{ga.gch}) is a pseudo-distance on $\CC$.
\end{proposition}

In fact, this is equivalent to the construction of corollary
\ref{inv.gch}. To see this, let as before  $G$ be the stabilizer of
$m_0$ for the action of $\CO$ on $\CC$. We have
\begin{proposition}
\label{cdis.gch}
If $\Ga$ satisfies C1), C2) and C3), then, for all $m_0\in\CC$, the
function  $d_\CO$ defined by
$$
d_\CO(o, o') = \Ga((o')^{-1}o, o^{-1}.m_0)
$$
is a distance on $\CO$ which is equivariant for the left action of
$G$. Conversely, given such a distance $d_\CO$, one build an effort
functional  $\Ga$ satisfying C1), C2), C3) b letting
$$\Ga(h, m) = d_\CO(oh, o)$$
where $o$ is any element of $\CO$ with the property $o.m = m_0$.
\end{proposition}
Here again, we leave the proposition to the reader.


\subsection{Infinitesimal approach}
\label{var}
The previous sections have demonstrated the interest of equivariant
distances on $\CO$. It may be, however, very difficult to come up with
an analytical expression for such a distance. Very often, fortunately
(in particular when  $\CC$ is a differential manifold and $G$ is a Lie
group) a left invarant Riemannian metric  can be designed, leading to a
left-invariant distance defined as a variational problem.

We need to build a left invariant metric on $\CO$. Recall that a
Riemannian metric on $\CO$ defines, for all $o\in\CO$, a dot product
$\scp{.}{.}_o$ on the tangent space, $T_o\CO$,  to $\CO$ at $o$, which
depends smoothly on $o$. With such a metric, one can define
the energy of a differentiable path  $\bfo$ in $\CO$ by
\begin{equation}
\label{ener.1}
E(\bfo) = \int_0^1 \left\|\der{\bfo}{t}\right\|^2 dt.
\end{equation}
The associated Riemannian distance on $\CO$ is
\begin{equation}
\label{geod.1}
d_\CO(o, o') = \inf\{ \sqrt{E(\bfo)}, \bfo(0) = o, \bfo(1) = o'\}.
\end{equation}

To obtain a left invariant distance, it suffices to ensure that the
metric has this property. For $h\in G$, we left $L_h$ denote the left
action of $h$ on $\CO$: $L_h: o \mapsto h.o$. Let $d_oL_h$ be its
differential at $o\in \CO$. It is a linear map, defined on 
$T_o\CO$, with values in $T_{h.o}\CO$ such that, for all
differentiable path $\bfo$ on $\CO$ such that $\bfo(0)=o$, we have
\begin{equation}
\label{left.trans}
\der{h\bfo}{t}_{|_{t=0}} = d_oL_h.\der{\bfo}{t}_{|{t=0}}.
\end{equation}
The left invariance of the metric is expressed by the identity, true
for all  $o\in \CO$, $A\in T_o\CO$ and $h\in G$,
\begin{equation}
\label{left.inv.1} 
\|A\|_o = \left\|d_oL_h.A\right\|_{ho}.
\end{equation}

When  $\CO = G\ti \CC$, condition
(\ref{left.inv.1}) implies that it suffices to define $\scp{.}{.}_o$
at elements  $o\in \CO$ of the form $o = (\id, m)$ with $m\in
\CC$. The metric at a generic point $(h,m)$ can then be computed, by
left invariance, from the metric at $(\id, h^{-1}m)$. Since the metric
at $(\id, m)$
can be interpreted as a way to attribute a cost to a deformation
$(\id, h(t)m)$ with $h(0) = \id$ and small $t$, defining it
corresponds to an analysis of the cost of an infinitesimal
perturbation of $m$ by an element of $G$.

In the next sections, we will apply this construction to several
situations, depending on the choice of the deformable objects and on
the group $G$.


\section{Invariant distances between landmarks}

\subsection{Introduction}
The purpose of this section is to present the construction provided by
Kendall \cite{ken84} on distances between landmarks, under the
infinitesimal point of view that we have just outlined. Here,
configurations of landmarks are considered up to similitude
(translation, rotation, scaling).

We only consider the 2D case, which is also the simplest. For a fixed
integer  $N >0$ let $\CP_N$ denote the set of configurations of 
$N$ points $(z_1, \ldots, z_N) \in (\mR^2)^{N}$. We assume that the
order in which the points are listed maters, which says that we
consider labelled landmarks. The set $\CP_N$ can therfore be ientified
to 
$\mR^{2N}$.

If the order between points is important, we will consider however,
that their global location in the 2D space is irrelevant.Two
configurations $(z_1, \ldots, z_N)$ and $(z'_1, \ldots, z'_N)$ will be
identified  if one can be deduced from the other by the composition
$g$ of a translation and a plane similitude, i.e., $z'_i = g.z_i$ for
all $i$. Objets of interest therefore are equivalent classes of
landmark configurations, which will be referred to as
$N$-shapes.

It will be convenient to identify the plane  $\mR^2$ with the set of
complex numbers $\mC$, a point $z =
(x,y)$ being represented as $x+iy$. A plane similitude composed with a
translation canthen be written under the form
$z \mapsto a.z + b$ with $a, b\in \mC$, $a\neq 0$.

For $Z = (z_1, \ldots, z_N)\in \CP_N$, we let $\bz$ be the center of inertia
$$
\bz = (z_1+\cdots+z_N)/N.$$
We also let  $\dps \|Z\|^2 = \sum_{i=1}^N |z_i-\bz|^2$.


\subsection{Space of planar $N$-shapes}

\subsubsection{Construction of a distance}
Let $\Sig_N$ be the quatient space of  $\CP_N$ by the equivalent
relation  $Z\,\CR\, Z'$ if there exists $a, b\in\mC$ such that
$Z'= aZ+b$. We denote $[Z]$ th equivalence class of 
$Z$. We shall define a distance between two equivalence classes $[Z]$
and $[Z']$. 

We can apply the results of the previous sections, which ensures that,
in order to define a distance in the quotient space, it suffices to
start with a distance on  $\CP_N$ which is equivariant for the
considered operations. So we need a  distance $D$ such that for all
$a, b\in\mC$ and all $W, Z\in \CP_N$,
$$
D(aZ+b, aW+b) = D(Z, W).
$$

To obtain such a distance, it suffices to definr a Riemannian metric
on $\CP_N = 
\mR^{2N}$ which is invariant by the action. We therefore must define,
for all  $Z\in \CP_N$, a norm $\|A\|_Z$ defined for all $A = (a_1,
\ldots, a_N)\in 
\mC^N$ such that for all $a, b\in \mC$:
$$
\|A\|_Z = \|a.A\|_{a.Z + b}.
$$
This formula implies that i suffices to define such a norm for  $Z$
such that $\|Z\| = 1$ and $\bz = 0$, since we have, for all
$Z$,
$$
\|A\|_Z = \|\frac{1}{\|Z\|}.A\|_{\frac{Z-\bz}{\|Z\|}}.
$$
Once the metric has been chosen, the distance $D(W, Z)$ is defined by
$$
D(W, Z)^2 = \inf \int_0^1 \norm{\der{\bfZ}{t}}^2_{\bfZ(t)} dt
$$
the infimum being taken over all paths  $\bfZ(.)$ such that $\bfZ(0)
= W$ and $\bfZ(1)  = Z$.

Le's make the following choice. When $\bz = 0$ and $\|Z\|=1$, we let
$$
\|A\|_Z = \sum_{i=1}^N |a_i|^2
$$
so that we need to minimize, among all paths between  $W$ and $Z$,
$$
\int_0^1
\frac{\sum_{i=1}^N|\dot{\bfZ_i}(t)|^2}{\sum_{i=1}^N|\bfZ_i(t)-\bar{\bfz}(t)|^2}
dt.
$$
Denote $\bfv_i(t) = (\bfZ_i(t) - \bfz(t))/\|\bfZ(t)\|$ and $\rho(t) =
\|\bfZ(t)\|$. The path $\bfZ(.)$ is uniquely $(\bfv(.), \rho(.),
\bz(.))$. Moreover, we have
$$
\der{\bfZ_i}{t} = \der{\bfz}{t} + \rho \der{\bfv}{t} + \bfv \der{\rho}{t}
$$
so that we need to minimize
$$
\int_0^1 \sum_{i=1}^N
\left|\frac{\der{\bfz}{t}}{\rho}+\frac{\der{\rho}{t}}{\rho}.\bfv_i + 
\der{\bfv_i}{t} \right|^2 dt.
$$
This is equal (using $\sum_i \bfv_i = 0$ and $\sum_i|\bfv_i|^2 =
1$, together with the differential on these expressions) to
$$
N\int_0^1 \left(\frac{\der{\bfz}{t}}{\rho}\right)^2
+ \int_0^1 \left(\frac{\der{\rho}{t}}{\rho}\right)^2 dt + \int_0^1
\sum_{i=1}^N \left|\der{\bfv_i}{t}\right|^2 dt\,.
$$
The boundary conditions directly deduced from those on  $W$ and $Z$.

The last term of this expression, which only depends on $\bfv$, can be
minimized explicitly, under the constraints 
$\sum_i \bfv_i = 0$ and $\sum_i|\bfv_i|^2 = 1$, which imply that
$\bfv$ varies on a $(N-1)$-dimensional sphere. The geodesic distance
is therefore given by the length of great circles, which yields the
expression of the minimum of the last term:
$ \arccos(\scp{\bfv(0)}{\bfv(1)})^2$.

We therefore have
\begin{eqnarray*}
D(W, Z)^2 &=& \inf \left(N\int_0^1 \left(\frac{\der{\bfz}{t}}{\rho}\right)^2
+ \int_0^1 \left(\frac{\der{\rho}{t}}{\rho}\right)^2 dt\right)\\
&& + \arccos\left(\scp{\frac{W-\bw}{\|W\|}}{\frac{Z-\bz}{\|Z\|}}\right)^2,
\end{eqnarray*}
where the first infimum is over functions $t\mapsto
(\bfz(t), \rho(t)) \in \mC \ti [0, +\infty[$, such that $\bfz(0) =
\bw$, $\bfz(1) = \bz(1)$, $\rho(0) = |W|$, $\rho(1) = |Z|$.

The induced distance on  $\Sig_N$ is then given by
$$
d([Z], [W]) = \inf\{D(Z, aW+b), a, b\in \mC\}.
$$
Since we can always select  $|a|$ and $b$ such that $|Z| =
|aW+b|$ and $\bz = a \bw + b$, we have
$$
d([Z], [W]) = \inf_{a}
\left[\arccos\left(\scp{\frac{a}{|a|}\frac{W-\bw}{\|W\|}}{\frac{Z-\bz}{\|Z\|}}\right)\right].
$$
Finally, optimizing this over the unit vector $a/\abs{a}$, we get
\encadre{
\begin{equation}
\label{dist.pol}
d([Z], [W]) =
\arccos\left|\scp{\frac{W-\bw}{\|W\|}}{\frac{Z-\bz}{\|Z\|}}\right|.
\end{equation}
}

Remark that the class $[Z]$ is characterized by $v = (Z-
\bz)/\|Z\|$, which varies in a complex sphere of dimension
$N-2$, denoted $\tS^{N-2}$, because it is the intersection of the
$(N-2)$-dimensional sphere and the 
complex hyperplane $\bz=0$. We have the equivalence
l'�quivalence
\begin{equation}
\label{carte}
[Z] = [Z'] \Leftrightarrow \exists \nu \in\mC: |\nu| = 1 \text{ and }
\Phi(Z) = \nu \Phi(Z')
\end{equation}

Denote by $S^{2N-3}$ the set of  $v = (v_1, \ldots, v_{N-1})\in
\mC^{N-1}$ such that 
$\sum_i |v_i|^2 = 1$ (this can be identified to a real sphere of
dimension  $2N-3$). The complex projective space, denoted $\mC
P^{N-2}$, is defined as the space $S^{2N-3}$ quotiented by the
equivalence relation $v\CR v'$ if and only if $\exists
\nu\in\mC$ such that $v'= \nu v$; in other words, $\mC P^{N-2}$
contains all sets
$$
S^1.v = \{\nu.v, \nu\in\mC, |\nu| = 1\|$$ when $v$ varies in
$S^{2N-3}$ (this is also called the {\em Hopf fibration} of spheres of
odd dimension). This set has a strcte of $(N-2)$ dimensional complex
manifold, which means that it can be covered by an atlas of open set
which are in bijection with open subsets of  $\mC^{N-2}$. It suffices
indeed to define  $\CO_i$ to be the set of $S^1. v
\in\mC P^{N-2}$ with $v_i\neq 0$, and then to set $\Psi(S^1 v) =
(v_1/v_i, \ldots, v_{i-1}/v_i, v_{i+1}/v_i, \ldots, v_{N-1}/v_i)
\in \mC^{N-2}$ for all $S^1 v\in \CO_i$. This provides $\mC
P^{N-2}$ with a structure of $C^\infty$ complex manifold. The
equivalence (\ref{carte}) therefore provides a representation of 
$\Sig_{N}$ as a $C^{\infty}$ complex manifold.

Let us be more precise with this identification \cite{ken84}. We
associate to $Z = (z_1, \ldots,
z_N)$ the family $(\ze_1, \ldots, \ze_{N-1})$ defined by
$$
\ze_i = (jz_{j+1} - (z_1+\cdots+z_j))/\sqrt{j^2+j}.
$$
One can verify that  $\sum_{i=1}^{N-1} |\ze_i|^2 = \|Z\|^2$. Denote
$F(Z)$ the element $S^1.(\ze/\|Z\|)$ in $\mC P^{N-2}$. One checks that
$F(Z)$ only depends on $[Z]$ and that $[Z]\mapsto F(Z)$ is a bijection
between  $\Sig_N$ and $\mC P^{N-2}$.


\subsubsection{The space of triangles}
Consider the case $N=3$, which coresponds to labelled triangles. For a
triangle $Z = (z_1, z_2, z_3)$, the previous function $F(Z)$ can be
written 
\begin{eqnarray*}
F(Z) &=& S^1.\left\{\frac{1}{\sqrt{|z_2-z_1|^2/2 + |2z_3 - z_1 -
z_2|^2/6}}\left[ \frac{z_2-z_1}{\sqrt{2}}, \frac{2z_3 - z_1 -
z_2}{\sqrt{6}} \right] \right\}\\
&=& S^1.[v_1, v_2].
\end{eqnarray*}
On the set $v_1 \neq 0$ (i.e., the set $z_1\neq z_2$) we have the
local chart
$$
Z \mapsto v_2/v_1 = \frac{1}{\sqrt{3}}\left(\frac{2z^3 - z_2 -
z_1}{z_2 -z_1}\right) \in\mC. 
$$
If we let $v_2/v_1 = \tan\frac{\th}{2} e^{i\phi}$, and $M(Z) =
(\sin\th\cos\phi, \sin\th\sin\psi, \cos\th) \in\mR^3$, we obtain a
correspondence between the  triangles and the unit sphere $S^2$. This
correspondence is isometric: the distance between two triangles  $[Z]$
and $[\tZ]$, which has been defined above by 
$$
d([Z], [Z']) = \arccos \left|\frac{\sum_{i=1}^3 z_i\tz_i^*}{\|Z\|.\|\tZ\|}\right|
$$
gives, after passing to coordinates $\th$ and $\phi$, exactly the
lenth of the great circle between the images  $M(Z)$ and $M(Z')$. We
therefore obtain a representation of labelled triangular shapes as
points on the sphere $S^2$, with the possibility of comparing them
using the standard metric on  $S^2$.


\subsection{Extension to plane curves}

It is interesting to analyze whether this distance has a limit when
the landmarks correspond to increasingly refined discretizations of
plane curves. Consider for this two plane curves $m$ and $\tm$,
parameterized on $[0,1]$. Discretize the interval  $[0,1]$ into $N+1$
points $t_k =  k/N, k=0,
\ldots, N$. Denote $Z_k = m(t_k)$ and $W_k = \tm(t_k)$. Let also 
$m^N$ and $\tm^N$ be the piecewise constant approximations of  $m$ and
$\tm$ given by $m^N(t) = Z_k$ (resp. $\tm^N(t) = W_k$)
on $]t_{k-1}, t_k]$.

We have $\dps \bz = (1/N)\sum_k m(t_k) \simeq \int_0^1 m(t) dt$. Similarly
$$
\|Z\|^2  \simeq N \left[\int_0^1 |m(t)|^2 dt - \left|\int_0^1
m(t)dt\right|^2\right]
$$
and
$$
\scp{Z}{W} \simeq N \left[\int_0^1 m(t)\tm(t)^* dt - \int_0^1
m(t)dt\int_0^1\tm(t)^*dt\right]
$$
which gives, for large  $N$,
$$
d(m, \tm) = \arccos\left\{\frac{\dps \left|\int_0^1 m(t)\tm(t)^* dt - \int_0^1
m(t)dt\int_0^1\tm(t)^*dt\right|}{\left[ V_m\right]^{1/2} \left[V_{\tilde m}\right]^{1/2}}\right\}
$$
with $V_m = 
\dps \int_0^1 |m(t)|^2 dt - \left|\int_0^1
m(t)dt\right|^2$ (and a similar definition for $V_{\tilde m}$).

We can get rid of the translation component by comparing the
differences $Z_k = m(t_k) - m(t_{k-1})$ instead of the points
$m(t_k)$. Doing so, we obtain the same formula, replacing  $m(t)$ by
$dm/dt$. Since 
$$
\int_0^1\der{m}{t}dt = m(1) - m(0),
$$
we have, if $m$ and $\tm$ are closed, which will be assumed for
simplification,
$$
d(m, \tm) = \arccos\left\{\frac{\dps \left|\int_0^1
\der{m}{t}\der{\tm}{t}^* dt\right|}{\left[\int_0^1
\left|\der{m}{t}\right|^2 dt \right]^{1/2} \left[\int_0^1
\left|\der{\tm}{t}\right|^2 dt \right]^{1/2}}\right\}.
$$
(If the cures are open, the obtanied formulae
coresponds to comparing them after application of the closing
operation  $m(t) \mapsto m(t) - t (m(1)-m(0))$.)

We now introduce a parameterization in this formula. Assume that the
curves have length one (which is no loss of generality, since the
distance is scale-invariant). We let  $s = \phi(t)$ for curve $m$ and $s
= \psi(t)$ for curve $\tm$, where $s$ is the arc-length. We then have
$\der{m}{t} = \dphi(t) 
\tau(\phi(t))$ where $\tau$ is the unit tangent (expressed in function
of the arc-length). To keep with complex notation, we let $\tau =
e^{i\th}$. We then have 
$$
\int_0^1 \left|\der{m}{t}\right|^2 dt = \int_0^1 \dphi(t)^2 dt
$$
and
$$
\int_0^1 \der{m}{t}\der{\tm}{t}^* dt = \int_0^1 \dphi(t)\dpsi(t)
e^{i(\th(\phi(t))-\tth(\psi(t)))} dt.
$$
This yields the formula
$$
d(m, \tm) = \arccos \left\{\frac{\left|\int_0^1 \dphi(t)\dpsi(t)
e^{i(\th(\phi(t))-\tth(\psi(t)))} dt\right|}
{\sqrt{\int_0^1 \dphi(t)^2 dt}\sqrt{\int_0^1 \dpsi(t)^2 dt}}\right\}.
$$
The interest in making the parameterization apparent is that this
offers the possibility for optimizing the distance with respect to it,
which relates to the optimal matching problems discussed in the
previous section.




\section{Distances and Diffeomorphisms}
\label{dist.diff}
\subsection{Invariant Distances between Diffeomorphisms}
\label{dist.inv.diff}

Our goal here is to build a right invariant distance between
diffeomorphisms. This study has its own interest, but will also be
useful when we will deal with  more general
actions involving diffeomorphisms. It will also allow us to retrieve
the construction that we have made previously for the groups of
diffeomorphisms.

We only make here a formal (non-rigorous) discussion. We consider a
group  $G$ of diffeomorphisms of $\Om$, and define (to fix the ideas)
the tangent space to $G$ at $g$ by the set of  $u:\Om\to\mR^d$ such
that $\id + tu\circ
g^{-1}\in G$ for small enough $t$. Since the ``product'' on $G$ is the
composition, $gh = g\circ h$, the right translation $R_g:
h\mapsto h\circ g$ is linear, and therefore equal to its differential:
for 
$u\in T_hG$:
$$
d_h R_g . u = u\circ g.
$$

A metric on  $G$ is right invariant if, for all  $g, h\in G$ and for
all $u\in T_hG$:
$$
\|d_h R_g u\|_{h\circ g} = \|u\|_h
$$
which yields, taking $g = h^{-1}$:
$$
\|u\|_h = \|u\circ h^{-1}\|_\id.
$$


This implies that the energy of a path $(t\mapsto \bfg(t, .))$ in 
$G$ must be defined by
$$
E(\bfg) = \int_0^1 \left\|\pdr{\bfg}{t}(t, \bfg^{-1}(t, .))\right\|^2_\id dt.
$$
Define 
$$\bfv(t, x) = \pdr{\bfg}{t}(t, \bfg^{-1}(t, x))$$
so that the energy can be written
$$
E(\bfg) = \int_0^1 \left\|\bfv(t,.)\right\|^2_\id dt
$$
with the identity
$$
\pdr{\bfg}{t}(t, x) = \bfv(t, \bfg(t, x)).
$$
This implies that $\bfg$ is the flow associated to the velocity field
$\bfv(t, .)$, which provides the construction given in chapter \ref{grp.diff}.

The construction of $G$ directly from flows of diffeomorphisms
therefore becomes natural, as well as the condition that  $\|.\|_\id$
is admissible, according to definition \ref{adm}. The right invariant
distance on  $G$ therefore is 
$$
d(g_0, g_1) = \inf\sqrt{\int_0^1 \left\|\bfv(t,.)\right\|_\id^2 dt}
$$
where the minimum is taken over all $\bfv$ such that, for all
$x\in\Om$, the solution of the ODE
$$
\der{y}{t} = \bfv(t,y)
$$
with initial conditions $y(0) = g_0(x)$ is such that $y(1) =
g_1(x)$. This provides an interpretation of the cost function defined
in section \ref{viscous}, chapter \ref{est.diff} as a squared distance
to the identity
$$
E(g)= d(\id, g)^2.
$$

An explicit computation of this distance is generally impossible, but
there exists an exception, in 1D. Take  $\Om = [0,1]$
and
$$
\|u\|^2_\id = \int_0^1 \left| u_x\right|^2 dx.
$$
This not is not admissible, because it cannot control the supremum
norm of  $u_x = \prt u/\prt x$. One can define anyway the energy of a
path of diffeomorphisms  $\bfg(t,.)$
by
$$
U(\bfg) = \int_0^1 \int_0^1 \left|\pdr{\ }{x}\left(\pdr{\bfg}{t}\circ \bfg(t,
.)\right)\right|^2dxdt.
$$
This gives, after expanding the derivative and making the change of
variables  $ x = \bfg(t, y)$:
$$
U(\bfg) = \int_0^1 \int_0^1 \left|\pddr{\bfg}{t}{x}\right|^2
\left|\pdr{\bfg}{x}\right|^{-1} dy dt.
$$
Define $q(t, y) = \sqrt{\der{\bfg}{x}(t,y)}$: we have
$$
U(\bfg) = 4 \int_0^1 \int_0^1 \left|\der{q}{t}\right|^2 dy dt
$$
which yields
$$
U(\bfg) = 4 \int_0^1\norm{\der{q}{t}(t, .)}_2^2 dt.
$$
If the problem was to minimize this energy under the constraints  $q(0, .) =
\sqrt{\der{g_0}{x}}$ and $q(1, .) =
\sqrt{\der{g_1}{x}}$, the solution $q$ would be given by the line segments
$$
q(t, x) = t q(1, x) + (1-t) q(0, x).
$$
There is however an additional constraint that comes from the fcat
that $q(t, .)$ must provide a homeomorphism of $[0,1]$ for all $t$,
which implies  $\bfg(t, 1) = 1$, or, in terms of $q$
$$
\|q(t, .)\|_2^2 = \int_0^1 q(t, x)^2 dx = 1.
$$

We therefore need to minimize the length of the path $q$ under the
constraint that it remains on a Hilbert ($\CL^2$) sphere. Like in
finite dimension, geodesics on Hilbert spheres are great circles. This
implies that the optimal  $q$ is given by
$$
q(t, .) = \frac{1}{\sin \al} (\sin(\al(1-t)) q_0 + \sin(\al t) q_1)
$$
with $\al = \arccos\scp{q_0}{q_1}_2$. The length of  the geodesic is
precisely given by  $\al$, which provides a closed-form expression of
the distance on  $G$
\encadre{
$$
d(g_0, g_1) = 2\arccos\int_0^1 \sqrt{\der{g_0}{x}\der{g_1}{x}} dx.
$$}

If one defines the cost of a diffeomorphism $g$ by its squared
distance to the identity, one gets
\encadre{
$$
E(g) = 2\arccos\int_0^1 \sqrt{\der{g_0}{x}} dx
$$}
\noindent which has been used in \cite{pst98}.


\subsection{Distance between deformable objects}
\subsubsection{Distance between Landmarks}

We here define a metric on the set $\Om^N$, composed with $N$-tuples
of labelled landmarks. Let $G$ be a group of diffeomorphisms, as built
in chapter  \ref{grp.diff}. We consider the space $\CO = G\ti
\Om^N$. The action of  $G$ on $\Om^N$ is
$$
g.(x_1, \ldots, x_N) = (g(x_1), \ldots, g(x_N))
$$
and the right action on $\CO$ is
$$
(h, x_1, \ldots, x_N)g = (h\circ g, g^{-1}(x_1), \ldots, g^{-1}(x_N)).
$$
The derivative of the right action therefore is
$$
dR_g(\xi, \al_1, \ldots, \al_N) = (\xi\circ g, Dg^{-1}(x_1)\al_1,
\ldots, Dg^{-1}(x_N)\al_N).
$$

We want to apply the construction given in section \ref{var}, and
therefore build a metric on $\CO$ which is right invariant by the
action of $G$. This metric must satisfy, for all  $(g, x_1, \ldots,
x_N)\in\CO$:
\begin{multline*}
\|(\xi, \al_1, \ldots, \al_N)\|_{g, x_1, \ldots, x_N} = \\
\|(\xi\circ
g^{-1}, Dg(x_1) \al_1, \ldots, Dg (x_N) \al_N)\|_{\id, g(x_1),
\ldots, g(x_N)}
\end{multline*}

Let a curve $\bfo(t) = (\bfg(t), \bfx_1(t), \ldots, \bfx_N(t))$ be
given in  $\CO$. Denote $y_i = \bfg(t,\bfx_i)$ and introduce 
$\bfv$ such that
$$
\pdr{\bfg}{t} = \bfv(t, \bfg).
$$

The previous formula provides the energy of the curve, which is
\begin{eqnarray*}
E(\bfo(.)) &=& \int_0^1 \norm{\der{\bfo}{t}}_{\bfo(t)}^2 dt\\
 &=& \int_0^1 \norm{\bfv,
\der{y_1}{t} - \bfv(t, y_1),\ldots, \der{y_N}{t} - \bfv(t, y_N)
}_{\id, y_1, \ldots, y_N}^2 dt\\
&=&\tE(\bfv, y)
\end{eqnarray*}
using the fact that 
$$
\der{y_i}{t} = \der{}{t} \bfg(t, x_i) = \bfv(t, y_i) +
D\bfg(x_i)\der{x_i}{t}.
$$

The geodesic distance associated to  $D$ on $\CO$ can be obtained by
minimizing this energy with fixed boundary conditions. These
conditions are, if the
compared objects are  $o_0$ et $o_1$, 
$$
o_0 = \bfo(0) = (\bfg(0)^{-1}, \bfg(0, y_1(0)), \ldots, \bfg(0,
y_N(0)))
$$
and
$$
o_1 = \bfo(1) = (\bfg(1)^{-1}, \bfg(1, y_1(1)), \ldots, \bfg(1,
y_N(1))).
$$

The projected distance between two landmark configurations is then
given by, denoting  $y^0 = (y^0_1, \ldots, y^0_N)$ and $y^1 = (y^1_1, \ldots,
y^1_N)$,
\begin{eqnarray*}
d(y^0, y^1) &=& \inf\defset{D(o_0, o_1): o_0 = (g^{-1}, g.y^0), o_1 = (h^{-1},
hy^1)}\\
&=& \inf\defset{\sqrt{E(\bfo)}: o(0) = (g^{-1}, g.y^0), o(1) = (h^{-1},
h.y^1)}\\
&=& \inf\defset{\sqrt{\tE(\bfv, y)}: y(0) = y^0, y(1) = y^1}.
\end{eqnarray*}
We can see that the last expression does not involve the
diffeomorphism in the boundary conditions: $\bfv$ becomes an
auxilliary variable for the minimization. This will allow us to define
a new form of interpolating splines.

For this, let's make the following choice for the norm at the identity
$$
\|\xi, \al_1, \ldots, \al_N\|^2_{\id, y_1, \ldots, y_N} = \|\xi\|^2 +
\la \sum_{i=1}^N |\al_i|^2,
$$
with an admissible norm for $\xi$, so that
\encadre{
$$
\tE(\bfv, y) = \int_0^1 \|\bfv\|^2 dt + \la\sum_{i=1}^N \left|
\der{y_i}{t} - \bfv(t, y_i)\right|^2
$$}

If the norm that is used for $\bfv$ comes from a dot product
associated to a kernel $K: (x, y)\mapsto K(x, y)$, with values in the
set of $d$ by $d$ matrices, such that, for all $a\in \mR^d$ and
$x\in\Om$, $a^T v(x) =  \scp{K(x,
.)a}{v}$. With the same arguments as in section \ref{joshi}, chapter
\ref{est.diff}, the optimal path $\bfv$ is such that there exist
vectors $a_1(t), \ldots, a_N(t)$ such that, for all  $t$:
$$
v(t, x) = \sum_{i=1}^N K(x, y_i(t))a_i(t).
$$
This implies that $\tE$ can be written
\begin{multline*}
\tE(a, y) = \sum_{i, j=1}^N \int_0^1 {}^ta_i(t) K(y_i(t), y_j(t))
a_j(t)dt \\
+ \la\sum_{i=1}^N \int_0^1 \left|\der{y_i}{t} - \sum_{j=1}^n
K(y_i(t), y_j(t)) a_j(t)\right|^2 dt.
\end{multline*}
The vectors $a_1, \ldots, a_N$ can be computed explicitly given the
trajectories $y_1, \ldots, y_N$, namely
$$
a = (S(y) + \la I)^{-1} \der{y}{t}.
$$
This provides an expression of the energy in terms of $y_1, \ldots,
y_N$ that can be minimized by gradient descent, for example. Figure
\ref{fig:splines0} provides some examples of deformations using this
method, called {\em geodesic splines} \cite{cy01}, with a comparison
to the linear techniques of section \ref{par:splines}, chapter \ref{est.diff}.

\begin{figure}
\begin{center}
\includegraphics[width=0.48\textwidth]{FIGURES/ini3.eps}\  
\includegraphics[width=0.48\textwidth]{FIGURES/path3.eps}\\ 
\includegraphics[width=0.48\textwidth]{FIGURES/diff3.eps}\  
\includegraphics[width=0.48\textwidth]{FIGURES/stat3.eps}
\end{center}
\caption{\label{fig:splines0} First Row: Random displacements of
16 points uniformly distributed on a grid (left). Optimal trajectories
for Geodesic Splines (right). Second Row: Grid deformation provided by
geodesic splines (left). Grid deformation provided by classic splines.} 
\end{figure}




\subsubsection{Distance between functions}
\label{cons.cons}
In this section, the object space  $\CC$ contains function  $m:
\Om\to\mR^d$. We use the same construction as before, the action of
$G$ on $\CC$ being here $g.m = m\circ g^{-1}$. We let again $\CO = G\ti \CC$.

The right action of $G$ on $\CO$ is $R_h(g,m) = (g\circ h, m\circ
h)$. If $o=(g, m)$, the tangent space to $\CO$ at $m$ contains pairs of
function  $(\xi, \eta)$ defined on $\Om$, and the diffrential of the
right translation at $o$ is given by $DR_h(o) (\xi, \eta) = (\xi \circ h,
\eta\circ h)$.
To define a right-invariant metric on, we must provide, for all $o$, a
norm
$$
(\xi. \eta)\mapsto \|\xi, \eta\|_o
$$
such that, for all $h$,
$$
\|\xi\circ h, \eta\circ h\|_{(g\circ h, m\circ h)} = \|\xi,
\eta\|_{(g, m)}.
$$
This implies that it suffices to define $\|\xi, \eta\|_{(\id, m)}$ for
all $m$ and to set
$$
\|\xi, \eta\|_{(g, m)} = \|\xi\circ g^{-1}, \eta\circ g^{-1}\|_{(\id, m\circ g^{-1})}.
$$

Given such a metric, we evaluate the energy of a path $\bfo(.)$ in
$\CO$, given as always by
$$
E(\bfo) = \int_0^1 \norm{\der{\bfo}{t}(t)}^2_{\bfo(t)}dt.
$$

Introduce, as before, $\bfv(t, .) = \pdr{\bfg}{t} \circ \bfg^{-1}(t,
.)$. Let also $\bfM(t, .) = \bfm(t, \bfg^{-1}(t, .))$. We have
$$
\pdr{\bfm}{t} (t, y) = \pdr{\bfM}{t}(t, g(t, y)) + \pdr{\bfM}{x}(t,
\bfg(t, y)).\pdr{\bfg}{t}(t, y)
$$
so that 
$$
E(\bfo) = F(\bfv, \bfM) := \int_0^1 \left\|\bfv(t, .),
\pdr{\bfM}{t}(t, .) + \pdr{\bfM}{x}(t,.).\bfv(t, .) \right\|_{\id, \bfM(t, .)}dt.
$$

Let $\bfg_\bfv$ denote the solution of
$$
\pdr{\bfg_\bfv}{t}(t, x) = \bfv(t, \bfg_\bfv(t, x))
$$
with $\bfg_\bfv(0, x) = x$ for all $x$. We obtain a distance
$d_\CO(o, o')$ by setting (with
$o=(g, m)$ and $o'= (g', \tm)$)
$$
d_\CO(o, o') = \inf \sqrt{F(\bfv, \bfM)}$$
where the infimum is taken over all paths $(\bfv, \bfM)$ such that
$\bfg_\bfv(1, g(x)) = g'(x)$, $\bfM(0, g(x)) = m(x)$ and $\bfM(1,
g(x)) = \tm(x)$ (this is simply the definition of the geodesic
distance on $\CO$ expressed in terms of the new variables  $\bfv$ and
$\bfM$).

The projected distance between two functions  $M$ and $M'$ is the
infimum of the previous expression, over all pairs  $o=(g, m)$ and
$o'= (g', \tm)$ such that
$m\circ g^{-1} = M$ and $\tm\circ {g'}^{-1} = M'$. By the right
equivariance of $d_\CO$,  it is also the infimum of 
$d_\CO[(\id, M), (g, \tm)]$ over all pairs $(g, \tm)$ such that
que $\tm\circ g^{-1} = M'$. Returning to the energy, it is therefore
the minimum of  $F(\bfv, \bfM)$ among all paths such that $\bfM(0,
x) = M(x)$ and $\bfM(1, \bfg_\bfv(1, x)) = \tm(x) = M(\bfg_\bfv(1,
x))$. It is therefore the  infimum of $F(\bfv, \bfM)$ over all paths
such that $\bfM(0, .) = M$ and $\bfM(1, .) = M'$. Here again, the
diffeomorphism does not appear in the boundary conditions.


\subsubsection{Application: comparison of plane curves}
\subsubsection*{Comparison based on orientation}

We first consider the issue of comparing plane curves on the basis of
the orientation of their tangents  \cite{you98}. If  $m$ is a plane
curve parameterized with arc-length, and $L$ is it length, we define
the normalized tangent  $T_m$ by
$$\map{T_m}{[0,1]}{S^1}{s}{\dm(Ls)}$$
where $S^1$ is the unit disc in $\mR^2$. The function $T_m$
characterizes $m$ up to translation and scaling, and the pair $(L,
T_m)$ characterizes $m$ up to translation.

We will identify, as we have done before, the sets $\mR^2$ and
complexe $\mC$. We consider a group of diffeomorphisms $G$ of $\Om
= (0,1)$, which acts on functions $T:[0,1]\to S^1$ par $gT =
T\circ g^{-1}$, and apply the construction of section
\ref{cons.cons} with
\begin{equation}
\label{n.tau}
\norm{\xi, \eta}^2_T= \norm{\xi, \eta}^2 = \int_0^1
\left[\left(\der{\xi}{x}\right)^2 + \eta^2\right] dx.
\end{equation}

The cost of the path $(\bfg(t, .), \bfT(t, .))$ is then
$$
U(\bfg, \bfT)= \int_0^1 \norm{\pdr{\bfg}{t}\circ \bfg^{-1},
\pdr{\bfT}{t}\circ \bfg^{-1}}^2 dt.
$$
Expand the derivative and make the change of  variables $x = \bfg(t,y)$
to write
$$
U(\bfg, \bfT) = \int_0^1 \int_0^1 \left[\left(\pddr{\bfg}{t}{x}\right)^2
\left(\pdr{\bfg}{x}\right)^{-1} +  \pdr{\bfg}{x}
\left|\pdr{\bfT}{t}\right|^2\right] dy dt
$$
Let $q(t, .)$ be differentiable in $t$ such that
$$q(t, x)^2 = \pdr{\bfg}{x}(t,x) \bfT(t,x)$$
($\bfT$ being considered as a unit complex number). Such a path always exists: for all  $x$, there exists (since $|\bfT(.,x)|=1$ a unique lifting  $\bfT(t, x) = \exp(i\th(t,x))$ with $\th(t,x)$
 differentiable in $t$ and $\th(0,x)\in[0, 2\pi[$. One can then define
\begin{equation}
\label{sqrt.q}
q(t, x) = \sqrt{\pdr{\bfg}{x}(t,x)} \exp(i\th(t,x)/2).
\end{equation}
This path is not uniquely defined, since one can always change the sign of $q(t,x)$, uniformly in $t$, for all $x$.

We have
$$
\pdr{q^2}{t}(t,x) = \pddr{\bfg}{t}{x}(t,x) \bfT(t,x)
+\pdr{\bfg}{x}(t,x) \pdr{\bfT(t,x)}{t}.
$$
Take the absolute value and use the fcat that $\bfT$ is orthogonal to
its derivative to obtain
$$
\abs{\pdr{q^2}{t}(t,x)}^2 = \left|\pddr{\bfg}{t}{x}(t,x)\right|^2
+\left(\pdr{\bfg}{x}(t,x)\right)^2
\left|\pdr{\bfT(t,x)}{t}\right|^2.
$$
Moreover
$$
\abs{\pdr{q^2}{t}(t,x)}^2 = 4|q(t,x)|^2\abs{\pdr{q}{t}(t,x)}^2 =
 4\pdr{\bfg}{x}(t,x)\abs{\pdr{q}{t}(t,x)}^2
$$
which implies
$$
|\pdr{q}{t}(t,x)|^2 = \frac{1}{4}\left[\left(\pddr{\bfg}{t}{x}\right)^2
\left(\pdr{\bfg}{x}\right)^{-1} +  \pdr{\bfg}{x}
\left|\pdr{\bfT}{t}\right|^2\right]
$$
so that
\begin{equation}
\label{eq:Uq}
U(\bfg, \bfT) = 4 \int_0^1 \int_0^1 \abs{\pdr{q}{t}(t,x)}^2 dx
dt.
\end{equation}


We now consider paths linking 
$(g_0, T_0)$ and $(g_1, T_1)$. We can already state that the  minimum
of $U(\bfg, \bfT)$ is larger than the minimum of the right-hand term of
\eqref{eq:Uq} over functions $q$ such that  $q(0,
.)^2 = \pdr{g_0}{x} T_0$ et $q(1, .)^2 = \pdr{g_1}{x} T_1$. We now
show that equality is true, provided that  $\pdr{g_0}{x}$ and
$\pdr{g_1}{x}$ are almost everywhere positive on  $[0,1]$.

First, note that, the computation the minimum of the right-hand
term of \eqref{eq:Uq} with fixed boundary conditions  $q_0$ and $q_1$
fix�es can be handled similarly to the distance for 1D diffeomorphisms
that we have computed at the end of section \ref{dist.inv.diff}, the
only difference here being that $q$ is a complex-valued
function. Since we still have the constraint that the integral of the
squared modulus of $q$ is 1, we are still computing geodesics on an infinite
dimensional sphere. The optimal path is therefore given by
$$
\bfq(t, x) = \frac{1}{\sin \al} (\sin(\al(1-t)) q_0(x) + \sin(\al t) q_1(x))
$$
with $\al = \arccos\scp{q_0}{q_1}_2$ and
$$
\scp{q_0}{q_1}_2 = \int_0^1 \Re(q_0(x) \bar{q}_1(x)) dx
$$
where $\bar{q}_1$ is the complex conjugate of $q_1$. For any choice of
$q_0$ and $q_1$ satisfying $q_0^2 =
\pdr{g_0}{x} T_0$ and $q_1^2 = \pdr{g_1}{x} T_1$, the optimal path
$\bfq$ uniquely satisfies the path  $(\bfg(t, .), \tau(t,
.)$ as long as $\bfq$ does not vanish.

Let us choose  $q_0$ and $q_1$ to ensure that $\al$ is minimal. This
is achieved by choosing, for all  $x$, $q_0(x)$ and $q_1(x)$ such that
$\Re(q_0(x) \bar{q}_1(x)) \geq 0$ (since $\arccos$ is
decreasing); we then have
\begin{eqnarray*}
|\bfq(t, x)|^2 &=& \frac{1}{\sin^2 \al} (\sin^2(\al(1-t)) |q_0(x)|^2 +
 \sin^2(\al t) |q_1(x)|^2 \\
 &&+ 2 \sin(\al(1-t))\sin(\al t) \Re(q_0(x)
 \bar{q}_1(x)))\\
&\geq & \frac{1}{\sin^2 \al} (\sin^2(\al(1-t)) |q_0(x)|^2 +
 \sin^2(\al t) |q_1(x)|^2).
\end{eqnarray*}
The last inequality shows that the function 
$\bfq(t,x)$ can vanish (for $t\in (0,1)$) only if both $q_0(x)$ et
$q_1(x)$ vanish. We are assuming that this can happen only when $x$
belongs to a set of Lebesgue's measure 0 in  $[0,1]$. On this set, the
function $t\mapsto\bfT(t,x)$ is undefined, but we can make an
arbitrary choice without changing the value of $\al$.

We have therefore showed that the distance between  $o_0 = (g_0, \tau_0=
e^{i\th_0})$ and $o_1  = (g_1,
\tau_1=e^{i\th_1})$ was given by 
$$
D(o_0, o_1) = 2 \arccos \int_0^1 \sqrt{\der{g_0}{x}\der{g_1}{x}}|\cos((\th_0(x) -
\th_1(x))/2)| dx
$$
since $|\cos((\th_0(x) -\th_1(x))/2)|$ is equal to  the real part of
the square root of  $T_0\bar{T}_1$ for which $\al$ is minimized. This
can also be written
$$
D(o_0, o_1) = 2\arccos \frac{\sqrt{2}}{2}\int_0^1
\sqrt{\der{g_0}{x}\der{g_1}{x}}\sqrt{1 +
\scp{\tau_0(x)}{\tau_1(x)}}dx .
$$
with $\scp{z}{z'} = \Re(z\bar{z}')$ for complex numbers $z$ and $z'$.

The resulting distance between two functions $T$ and $T'$ is,
according to  (\ref{inv.eq.gen})
\begin{eqnarray*}
d(T, T') &=& \inf\{D(o_0, o_1), o_0 = (g_0, T_0), o_1 = (g_1,
T_1), T_0\circ g_0^{-1} = T, T_1\circ g_1^{-1} = T'\}\\
&=& \inf\{D(o_0, o_1), o_0 = (\id, T), o_1 = (g,
T_1), T_0\circ g_0^{-1} = T, T_1 = T'\circ g^{-1} \}
\end{eqnarray*}
which finally yields
\encadre{
$$
d(\tau, \tau') = 2\arccos \sup_g \int_0^1 \sqrt{\der{g}{x}} |\cos((\th - \th'\circ
g)/2)| dx.
$$}

The existence (under some conditions on $\th$ and $\th'$) of an
optimal $g$ in the previous formula is proved in \cite{ty00}. Figure
\ref{fig:exp.1D} provides examples of curve matching that can be
obtained using this method.

\begin{figure}
\begin{center}
\includegraphics[width=5cm]{FIGURES/curveMatching1.ps}
\includegraphics[width=5cm]{FIGURES/curveMatching2.ps}\\
\includegraphics[width=4cm]{FIGURES/curveMatching3.ps}
\includegraphics[width=6cm]{FIGURES/curveMatching4.ps}\\
\end{center}
\caption{\label{fig:exp.1D} Diffeomorphic matching between curves. In
each image, the geodesic is computed between the outer curve and the
inner curve. Scaling is for visualization only (the computation is
scale invariant).}
\end{figure}

\subsubsection*{Infinitesimal analysis}
We now provide an infinitesimal justification for the choice of the
metric  (\ref{n.tau}), which also offers alternative definitions.

We therefore consider infinitesimal variations of plane curves. Let
$m$ be a parameterized curve, defined on  $[0, 1]$, such that, for all  $x$
$$\left|\der{m}{x}\right|\equiv L(m)$$
where $L(m)$ is the length of $m$ (this corresponds to the arc-length
normalized by the length). Let's consider a small perturbation along
$m$, denoted  $x \mapsto \ep
\De(x)\in\mR^2$. We assume that $\ep$ is small
and only keep terms of order one in the following computation. The
deformed curve is $x\mapsto m(x) + \ep \De(x)$, that we parameterize by
normalized arc-length to obtain a new curve $\tm$, still defined on
$[0,1]$. We denote by $g$ the diffeomorphism providing the change of
parameter, so that, for all  $x\in [0,1]$,
$$
\tm \circ g(x) = m(x) + \ep \De(x).
$$
Let $\xi(x) = g(x) - 1$. Let $T$, $N$ and $\ka$ be the tangent, normal
and Euclidean curvature of $m$. We have (denoting derivatives as subscripts)
$$
m_x = L(m) T,
$$
$$
m_{mx} = L(m)  T_x = L(m)^2 \ka N.
$$
Similarly, let  $\tT$, $\tN$, $\tka$, be the tangent, normal and
curvature of $\tm$; let finally $L = L(m)$ and $\tL = L(\tm)$. At
firest order, we have  
\begin{eqnarray*}
\ep \De_x &=& (1+\xi_x) \tm_x \circ g - m_x\\
&=& \tL \xi_x \tT\circ g + \tL\tT\circ g - L T\\
&=& L\xi_x T + (\tL-L)T  + L(\tT\circ g - T) + o(\ep).
\end{eqnarray*}
We now define an infinitesimal deformation energy in terms of $\De$
and compute its expression in terms of $g$ and the geoetric
characteristics of the curves. Denote 
$$
E = \int_0^L\left|V_x\right|^2 = \frac{\ep^2}{L} \int_0^1
\left|\De_x\right|^2dx
$$
where $V$ and $\De$ are related by $V(x) = \ep \De(x/L)$. At first
order in $\ep$, $\tT\circ g-
T$ is perpendicular to $T$, which implies 
$$
E = \frac{(\tL-L)^2}{L} + L \int_0^1\left|\xi_x\right|^2
dx + L \int_0^1|\tT\circ g - T|^2dx + o(\ep^2).
$$
Introduce the group $G$ of diffeomorphisms of  $[0,1]$, and consider
the set $\CO$ consisting of pairs $o = (L, T)$, with 
$L>0$ and $T: [0, 1]\mapsto \mR^2$ such that $|T(x)| = 1$ for all
$x$. There exists a unique curve  $m$ (modulo translation) with length
$L$ and having $T$ as a unit tangent. Denote
$$
\norm{\xi, \de_L, \de_T}_{L, \tau}^2 = L \int_0^1 \left|\xi_x\right|^2 dx
\frac{\de_L^2}{L} +
L\int_0^1|\de_T|^2dx.
$$
Without the length component, this is exactly the metric introduced in
the previous section. An explicit computation of the associated
geodesic distance is still possible. Let   $o_0 = (g_0,
L_0, T_0)$
and $o_1 = (g_1, L_1, T_1)$. We obtain \cite{you98}
$$
D(o_0, o_1) = L_0 + L_1 - \sqrt{2}L_0L_1 \int_0^1
\sqrt{\der{g_0}{x}\der{g_1}{x}}\sqrt{1 + \scp{T_0(x)}{T_1(x)}}dx
\ .
$$

\bigskip
Let us consider higher derivatives, which will allow us to also deal
with rotation invariance. We have
\begin{eqnarray*}
\ep \De_{xx} &=& L^2 \xi_x \ka N +
L\xi_{xx} T + L(\tL-L)\ka N + L\tL (1+\xi_x)\tka\circ g
N\circ g - L^2 \ka N +_ o(\ep)\\
& = & L\xi_{xx} T + 2L^2 \xi_x\ka N
+2L(\tL-L)\ka N + L^2(\tka\circ g - \ka)N + L^2\ka(\tN\circ g -
N) + o(\ep)
\end{eqnarray*}
At order 1, 
$L(\tN\circ g - N)$ is perpendicular to $N$, and can be computed by
applying a $\pi/2$-rotation to the vector
$L(\ttau\circ g - \tau)$, the latter being the normal component of the
first derivative of $\ep \De$. We can therefore write
$L^2(\tN \circ g - N) = -L\ep \scp{\De_x}{N} T$, which yields
$$
\ep \De_{xx} + L\ep\ka\scp{\De_x}{N} T = L\xi_{xx} T + 2L^2 \xi_x\ka N
+2L(\tL-L)\ka N + L^2(\tka\circ g - \ka) N. 
$$

This computation suggests the definition of the following operator
along $m$:
$$
E = \frac{\ep^2}{L^3} \int_0^1 \left|\De_{xx} +
L\ka\scp{\De_x}{N} T \right|^2 dx,
$$
which is invariant by (infinitesimal) rotation. (The factor $L^{-3}$
comes from an analysis using  $V(x) = \De(x/L)$.) This energy is equal
to 
$$
E = \frac{1}{L}\int_0^1 \left|\dder{\xi}{x}\right|^2dx +
L\int_0^1 \left|2\der{\xi}{x}\ka
+2\frac{\tL-L}{L}\ka + \tka\circ g - \ka\right|^2dx + o(\ep^2).
$$

This leads us to considering curve representation under the form $m =
(L, \ka)$, $L$ being the length, and  $\ka$ the curvature
parameterized by the normalized arc-length. The previous analysis
provides a natural candidate for the norm  $\norm{.}_{L,
\ka}$. Applying our general theory (section \ref{cons.cons}), this
defines the energy of  a path in $G\ti\CC$. Introduce the variable 
$\bfv$ as before, and the evolving object,  $(\bfL,
\bfK)$, which depend on time $t$. The energy of a path can then be
written as
$$
\int_0^1\int_0^1 \frac{1}{\bfL}\left|\dder{\bfv}{x}\right|^2dxdt +
\int_0^1\int_0^1 \bfL \left|2\der{\bfv}{x}\bfK
+2\frac{1}{\bfL}\bfK\der{\bfL}{t} + \pdr{\bfK}{t} +
\bfv\pdr{\bfK}{x}\right|^2dxdt.
$$




\subsubsection{Application: images matching}

We now provide a simple application of the formulas of section
\ref{cons.cons} for image comparison. The set $\Om$ is here a square
(or a rectangle) in $\mR^2$, for example $[0,1]^2$, and objects are
images, considered as functions defined on  $\Om$ and taking values in
$\mR$. We let
$$
\|\xi, \eta\|_{\id, m} = \int_\Om \scp{L\xi}{\xi}^2dx + \la \int_{\Om}
|\eta|^2dx
$$
where $L$ is a strictly monotonous admissible operator.

From the general formula, the energy of a path is given by
\begin{equation}
\label{eq:img.sign}
\tE(\bfv, \bfM) = \int_0^1\int_\Om \scp{L\bfv}{\bfv}^2 dx dt + \la \int_0^1
\int_\Om \left|\pdr{\bfM}{t} + \scp{\nabla \bfM}{\bfv}\right|^2dx dt.
\end{equation}

Figures \ref{fig:quadrature} and \ref{fig:nova1} provide two examples
of images matched using this energy. The numerical scheme is described
in the next section.



\begin{figure}
\begin{center}
\includegraphics[width=2.5cm]{FIGURES/quadra-newI01.eps}
\includegraphics[width=2.5cm]{FIGURES/quadra-newI03.eps}
\includegraphics[width=2.5cm]{FIGURES/quadra-newI06.eps}
\includegraphics[width=2.5cm]{FIGURES/quadra-newI11.eps}\\
\includegraphics[width=2.5cm]{FIGURES/quadra-grid01.eps}
\includegraphics[width=2.5cm]{FIGURES/quadra-grid03.eps}
\includegraphics[width=2.5cm]{FIGURES/quadra-grid06.eps}
\includegraphics[width=2.5cm]{FIGURES/quadra-grid11.eps}
\end{center}
\caption{\label{fig:quadrature} Estimation of a geodesic between a
disc and a square. First Row: Image evolution in time. Second Row:
evolution of the diffeomorphism.}
\end{figure}

\begin{figure}
\begin{center}
\includegraphics[width=2.5cm]{FIGURES/nova1-newI01.eps}
\includegraphics[width=2.5cm]{FIGURES/nova1-newI03.eps}
\includegraphics[width=2.5cm]{FIGURES/nova1-newI06.eps}
\includegraphics[width=2.5cm]{FIGURES/nova1-newI11.eps}\\
\includegraphics[width=2.5cm]{FIGURES/nova1-grid01.eps}
\includegraphics[width=2.5cm]{FIGURES/nova1-grid03.eps}
\includegraphics[width=2.5cm]{FIGURES/nova1-grid06.eps}
\includegraphics[width=2.5cm]{FIGURES/nova1-grid11.eps}
\end{center}
\caption{\label{fig:nova1} Geodesic between an empty image and a disc.
First Row: Image evolution in time. Second Row:
evolution of the diffeomorphism.}
\end{figure}

\subsubsection{Discretization of \eqref{eq:img.sign}}

Some care is required for the discretization of 
\eqref{eq:img.sign}. This is mainly due to the ``hyperbolic'' term
$\pdr{\bfM}{t} + \scp{\nabla \bfM}{\bfv}$, which represent the total
derivative of $\bfM$ along the flow. The stablest results have been
obtained using the following scheme.

Introduce the auxilliary variable $w = L^{1/2}v$, so that $v =
K^{1/2}w$ where $K= L^{-1}$ (the numerical implementation would in
fact define $\tK = K^{1/2}$ and let $K = \tK^2$). The following
discretized energy has been used in the experiments of Figures
\ref{fig:quadrature} and \ref{fig:nova1}. 
$$
E = \sum_{t=1}^T \sum_{x\in \tilde\Om} \abs{w(x)}^2 + \la\de_t^{-2} 
\sum_{t=1}^T \sum_{x\in \tilde\Om} \abs{M(t+1, x + \de_t (K^{1/2}w)(t, x))
- M(t, x)}^2
$$
Where $x$ and $t$ now are discrete variables, $\tilde\Om$ a
discrete grid on $\Om$ and $\de_t$ the time discretization step.



